% `mai` is the English Master of Science in Artificial Intelligence programme.
% `openany` prevents blank verso pages before chapters while retaining two-sided layout.
\documentclass[master=mai,openany]{kulemt}
\raggedbottom
\setup{
  title={Towards Safer Autonomous Navigation: Neurosymbolic Multi-agent RL in LogiCity Simulator},
  author={Yusuf Taskiran},
  promotor={Prof.\,dr.\,ir. \ Luc de Raedt \and Prof.\,dr.\,ir. Giuseppe Marra},
  assessor={Bart Bogaerts \and Renaud Detry},
  assistant={Ir.\ Ying Jiao \and Ir. Bahare Salmani}}
% Remove the "%" on the next line for generating the cover page
% \setup{coverpageonly}
% Remove the "%" before the next "\setup" to generate only the first pages
% (e.g., if you are a Word user).
%\setup{frontpagesonly}

% If you want to include other LaTeX packages, do it here. 

% Finally the hyperref package is used for pdf files.
% This can be commented out for printed versions.
\usepackage[pdfusetitle,colorlinks,plainpages=false]{hyperref}
\usepackage[table]{xcolor} % 'table'optie is handig voor cell backgrounds
\usepackage{xltabular}    % combinatie van longtable + tabularx
\usepackage[table]{xcolor}
\usepackage{enumitem}
\usepackage{xltabular,enumitem} 
\usepackage[table]{xcolor} 
\usepackage{pdflscape}
\usepackage{float}
\usepackage{subcaption}
\usepackage{amssymb} % for \checkmark
\usepackage{listings}
\usepackage{xcolor}
%%%%%%%
% The lipsum package is used to generate random text.
% You never need this in a real master's thesis text!
\IfFileExists{lipsum.sty}%
 {\usepackage{lipsum}\SetLipsumDefault{11-13}}%
 {\newcommand{\lipsum}[1][11-13]{\par And some text: lipsum ##1.\par}}
%%%%%%%
\usepackage{longtable}
%%%%%%%
\usepackage{amsmath}

\usepackage{listings}

\lstdefinelanguage{ProbLog}{
    morekeywords={query},
    sensitive=true,
    morecomment=[l]{\%},
    morestring=[b]"
}

\usepackage{listings}
\usepackage{xcolor}

\definecolor{codegreen}{rgb}{0,0.6,0}
\definecolor{codegray}{rgb}{0.5,0.5,0.5}
\definecolor{codepurple}{rgb}{0.58,0,0.82}
\definecolor{backcolour}{rgb}{0.95,0.95,0.92}

\lstdefinestyle{mystyle}{
    backgroundcolor=\color{backcolour},
    commentstyle=\color{codegreen},
    keywordstyle=\color{magenta},
    numberstyle=\tiny\color{codegray},
    stringstyle=\color{codepurple},
    basicstyle=\ttfamily\footnotesize,
    breakatwhitespace=false,
    breaklines=true,
    captionpos=b,
    keepspaces=true,
    numbers=left,
    numbersep=5pt,
    showspaces=false,
    showstringspaces=false,
    showtabs=false,
    tabsize=2
}

\lstset{style=mystyle}

\lstset{
    basicstyle=\ttfamily\small,
    columns=fullflexible,
    keepspaces=true,
    breaklines=true,
    frame=single,
    captionpos=b
}

\begin{document}

\begin{preface}
This is my master thesis to finish my master in Artificial Intelligence. I want to take a moment to acknowledge to thank everyone who has kept me busy.
\\
\\
I would like to thank my promotors Prof.dr.ir. Luc de Raedt and Prof.dr.ir. Giuseppe Marra, my daily advisors Ir. Ying Jiao and Ir. Bahare Salmani, and the entire jury.
\\
\\
I also wish to thank my parents for supporting me to finish my master and master thesis. 
\end{preface}

\tableofcontents*

\begin{abstract}
Multi-agent autonomous driving requires several vehicles to make safe, interdependent decisions while sharing road space, traffic rules, and limited local information. In such settings, improving the safety of an individual vehicle is insufficient: the system must also coordinate progress so that all vehicles can complete their routes without collisions or deadlock. This thesis investigates Probabilistic Logic Shielding (PLS) as a multi-agent reinforcement learning approach for safe task completion and coordination in the LogiCity Safe Path Following task. Each RL-controlled car receives a local, confidence-valued symbolic observation, and its probabilistic logic shield uses traffic rules to reweight a PPO policy towards safer speed-control actions. The shielded policy is used for both action selection and learning through a Probabilistic Logic Policy Gradient objective \cite{yang2023safereinforcementlearningprobabilistic}.
\\
\\
The evaluation uses compact intersection-centred scenarios containing one to three simultaneously controlled RL cars and rule-based pedestrians. It assesses scene-level trajectory success rate (TSR), failure rate, and timeout rate, distinguishing safe joint completion from collisions and coordination-induced non-completion. The experiments compare unshielded PPO with decentralized PLS as the number of RL cars increases; examine conservative, aggressive, and mixed decentralized shield configurations; analyse centralized shielding on representative decentralized deadlock cases; and compare shared-policy and separate-policy training for homogeneous agents.
\\
\\
Across all evaluated multi-agent configurations, conservative decentralized PLS improves TSR and reduces failures relative to unshielded PPO. However, as the number of interacting cars grows, some avoided failures become timeouts, revealing a trade-off between local caution and collective progress. A mixed two-car shield assignment achieves the strongest two-agent result, indicating that breaking behavioural symmetry can improve coordination. In selected deadlock cases, centralized shielding resolves mutual waiting by coordinating an ordered intersection passage. Shared-policy training also outperforms separate-policy training in the three-car setting. These findings show that PLS is a useful mechanism for multi-agent urban navigation, but that its value must be judged by its effect on collective safety, scene-level completion, and coordination rather than on individual failure reduction alone.
\end{abstract}

% A list of figures and tables is optional
%\listoffigures
%\listoftables
% If there are only a few figures and tables, it is better to use the following:
\listoffiguresandtables
% The list of symbols is also optional.
% This list must be created manually, e.g. as follows:
\chapter{List of abbreviations and symbols}
\section*{abbreviations}
\begin{flushleft}
  \renewcommand{\arraystretch}{1.1}
  \begin{tabularx}{\textwidth}{@{}p{12mm}X@{}}
    AD   & Autonomous Driving \\
    AI   & Artificial Intelligence \\
    CTDE   & Centralized Training with Decentralized Execution \\
    DRL   & Deep Reinforcement Learning \\
    DTDE   & Decentralized Training and Execution \\
    FOL   & First Order Logic \\
    FOV   & Field of View \\
    LiDAR   & Light Detection And Ranging \\
    MARL   & Multi Agent Reinforcement Learning \\
    MDP   & Markov Decision Problem \\
    MER   & Mean Episodic Reward \\
    MAS   & Multi Agent Systems \\
    PLPG   &  Probabilistic Logic Policy Gradient \\
    PLS   & Probabilistic Logic Shielding \\
    POMDP   & Partially Observable MDP \\
    PPO  & Proximal Policy Optimization \\
    RL   & Reinforcement Learning \\
  \end{tabularx}
\end{flushleft}
\section*{Symbols}
\begin{flushleft}
  \renewcommand{\arraystretch}{1.1}
  \begin{tabularx}{\textwidth}{@{}p{28mm}X@{}}

    $\pi_{\theta}(a \mid o)$ 
      & Base PPO policy probability of selecting action $a$ given observation $o$ \\

    $\pi_{\theta}^{+}(a \mid o)$ 
      & Shielded PPO policy after probabilistic safety reweighting \\

    $o$ 
      & Current symbolic observation/state \\

    $a$ 
      & Discrete action selected by the agent \\

    $L_{\mathrm{policy}}$ 
      & PPO clipped policy loss \\

    $L_{\mathrm{value}}$ 
      & PPO value-function loss \\

    $L_{\mathrm{entropy}}$ 
      & PPO entropy regularization term \\

    $L_{\mathrm{safety}}$ 
      & Safety regularization loss induced by probabilistic logic shielding \\

    $\alpha$ 
      & Safety regularization coefficient \\

    $c_v$ 
      & Weight coefficient for the value-function loss \\

    $c_e$ 
      & Weight coefficient for the entropy term \\

    $\log(\cdot)$ 
      & Natural logarithm used in the safety loss definition \\

    $\sum_{a}$ 
      & Summation over all possible actions \\

    $\theta$ 
      & Parameter vector of the PPO policy network \\

    $\mid$ 
      & Conditioning operator, read as ``given'' \\

  \end{tabularx}
\end{flushleft}

% Now comes the main text
\mainmatter

\chapter{Introduction}
\label{cha:intro}

Autonomous driving is a safety-critical sequential decision problem and a natural application of reinforcement learning (RL) \cite{kiran2021deepreinforcementlearningautonomous,isele2019safereinforcementlearningautonomous}. In multi-agent traffic, however, the consequences of one vehicle's decision depend on the behavior of nearby vehicles and pedestrians. Safe RL therefore needs to balance progress with traffic safety rather than optimize either in isolation.
\\
\\
This thesis studies Probabilistic Logic Shielding (PLS) in LogiCity, a logic-driven urban navigation environment with symbolic traffic predicates and rule-based safety conditions \cite{li2025logicityadvancingneurosymbolicai,yang2023safereinforcementlearningprobabilistic}. The implemented shield converts local symbolic traffic observations into graded action-safety scores and reweights the RL policy before action selection. This is well suited to relations such as intersection occupancy, priority, and collision proximity, for which a purely geometric or binary safety treatment is often insufficient.
\\
\\
The practical objective is \emph{safe task completion}: a vehicle must respect traffic rules while still reaching its destination within the episode horizon. This distinction is central in multi-agent traffic, where a shield may reduce unsafe failures but create excessive yielding and timeouts, or fail to resolve a shared intersection despite avoiding a collision.
Against this background, the main question of this thesis is:

\begin{quote}
\textbf{How does probabilistic logic shielding affect safe task completion and coordination in multi-agent LogiCity Safe Path Following?}
\end{quote}

\noindent We study this question through four research questions:

\begin{enumerate}
    \item How does PLS compare with unshielded RL as the number of simultaneously active RL-controlled cars increases?
    \item How do different decentralized shield styles change the trade-off between failure and timeout in multi-agent traffic interaction?
    \item Can centralized shielding resolve coordination failures that remain difficult under decentralized per-agent shielding?
    \item How does the choice between shared-policy and separate-policy multi-agent training affect performance and behavior under PLS?
\end{enumerate}

\noindent The experiments use a decentralized shared-policy baseline: multiple RL-controlled cars act simultaneously from local symbolic observations while sharing policy parameters. They compare unshielded RL with PLS as the number of controllable cars increases, evaluate conservative and aggressive decentralized shield variants, examine targeted centralized coordination in deadlock-prone interactions, and compare shared- and separate-policy training. Trajectory success rate is the primary outcome, interpreted alongside failure, timeout, intervention, and action-distribution measures.
\\
\\
We make three main contributions. First, we extend the study of probabilistic logic shielding to a multi-agent LogiCity SPF setting through a decentralized shield that combines local symbolic traffic observations with probabilistic action reweighting. Second, we evaluate PLS through safe task completion, distinguishing successful progress from unsafe failure and overly conservative timeout. Third, we systematically examine how PLS behaves under increasing multi-agent interaction and under key design choices in shielding, coordination, and policy organization.


\chapter{Background}
\label{cha:2}

\section{Autonomous Driving}

Autonomous driving (AD) systems are commonly described as a pipeline from sensing to control: sensor, perception, scene representation, planning/decision making, and control \cite{kiran2021deepreinforcementlearningautonomous}. Raw inputs from cameras, radar, LiDAR, and related sensors are transformed into a structured view of the local environment, from which the vehicle plans and executes actions.
\\
\\
Within this pipeline, decision making is especially challenging. Driving cannot be handled well by purely supervised learning because actions influence future states and interactions under uncertainty \cite{kiran2021deepreinforcementlearningautonomous}. This naturally motivates reinforcement learning (RL), where the vehicle learns a policy through sequential interaction with its environment.
\\
\\
Traffic is also fundamentally interactive. The ego vehicle shares the road with other vehicles and pedestrians, so its safe and effective behavior depends not only on its own actions but also on those of others \cite{zhang2024multiagentreinforcementlearningautonomous}. As a result, many driving situations, such as merging and intersection handling, are inherently multi-agent rather than purely single-agent planning problems.

\section{Multi-agent autonomous driving}

Urban traffic is a natural multi-agent system: autonomous vehicles interact with other vehicles, pedestrians, cyclists, and infrastructure. The survey by Zhang et al. argues that single-agent RL can lead to self-centered or overly aggressive behavior because it optimizes only the ego reward, whereas multi-agent formulations model coordination and mutual influence more directly \cite{zhang2024multiagentreinforcementlearningautonomous}.
\\
\\
Such environments are often formalized as Decentralized Partially Observable MDPs (Dec-POMDPs), where each agent has only local information and must act under uncertainty \cite{zhang2024multiagentreinforcementlearningautonomous}. Depending on the setting, interactions may be cooperative, competitive, or mixed. The same survey also highlights major training paradigms such as Centralized Training with Decentralized Execution (CTDE) and fully Decentralized Training and Execution (DTDE), both of which are relevant to traffic coordination.

\section{Deep Reinforcement Learning}

Autonomous driving can be modeled as sequential decision making in an MDP-like setting, where an agent repeatedly observes its environment and selects actions according to a policy. Deep reinforcement learning (DRL) combines this framework with deep neural networks, allowing policies to be learned from complex or high-dimensional inputs.
\\
\\
DRL has been applied to lane keeping, lane changing, merging, and intersection handling, using methods such as DQN variants, actor--critic algorithms including DDPG, A3C, and PPO, and hybrid RL-control architectures \cite{kiran2021deepreinforcementlearningautonomous}.
\\
\\
However, DRL alone is insufficient in safety-critical domains. Unsafe exploration and rare edge cases can produce unacceptable failures, so learning must be complemented by explicit safety constraints, prior knowledge, or formal guarantees \cite{isele2019safereinforcementlearningautonomous}. This motivates safe reinforcement learning.

\section{Safe Reinforcement Learning}

The core problem in applying RL to autonomous driving is the lack of safety guarantees. Standard RL learns through trial and error, which can lead to unsafe actions and even collisions during training or deployment. In real traffic, such behavior is unacceptable. Safe reinforcement learning (Safe RL) therefore seeks to optimize performance while respecting explicit safety constraints.

\subsection{Formalization of safety in reinforcement learning}

A widely used formal framework for Safe RL is the constrained Markov decision process (CMDP), where the agent maximizes expected return while keeping one or more cost functions below fixed thresholds \cite{gu2024reviewsafereinforcementlearning}. These cost functions represent safety violations and make the safety-performance trade-off explicit.
\\
\\
For autonomous driving, however, CMDPs have clear limitations. Many safety requirements are instantaneous rather than average-case, while traffic also involves partial observability and non-stationary multi-agent interaction. As a result, expected-cost constraints alone rarely provide strong or interpretable safety guarantees in realistic driving settings \cite{zhang2024multiagentreinforcementlearningautonomous}.

\subsection{Categories of Safe RL methods}

Following the review by Gu et al. \cite{gu2024reviewsafereinforcementlearning}, Safe RL methods can be grouped by how safety is enforced:

\paragraph{Reward-based and penalty-based methods.}
These methods add penalties for unsafe behavior directly to the reward. They are simple and widely used, but provide weak guarantees and depend strongly on reward design, which remains difficult in autonomous driving \cite{kiran2021deepreinforcementlearningautonomous}.

\paragraph{Constrained optimization methods.}
These methods incorporate safety constraints explicitly, often through primal-dual or Lagrangian techniques. They have a stronger theoretical basis than reward shaping, but typically enforce safety only in expectation, so violations may still occur during exploration.

\paragraph{Control-theoretic methods.}
These methods use tools such as Lyapunov functions, model predictive control, or control barrier functions to restrict actions to a safe set. They can provide strong guarantees, but often require accurate dynamics models and become computationally expensive in multi-agent traffic.

\paragraph{Formal methods--based methods.}
These methods use logic and verification to prove that unsafe states cannot be reached. They are attractive because safety is specified explicitly rather than learned implicitly, but they struggle with continuous state spaces, high-dimensional dynamics, and multi-agent interaction.

\paragraph{Gaussian Process--based methods.}
These methods model uncertainty in dynamics or costs and act conservatively in uncertain regions. They are useful for uncertainty-aware safety, but they also face scalability limitations.

\paragraph{Safety-layer and intervention-based methods.}
These methods enforce safety at runtime by filtering or correcting actions proposed by the learned policy. Shielding is a prominent example of this family.

\section{Shielding}

Shielding is a Safe RL technique that adds a formally synthesized runtime safety layer around a learned agent \cite{alshiekh2017safereinforcementlearningshielding}. Safety requirements such as collision avoidance, lane keeping, or traffic-rule obedience are encoded as logical specifications, and the shield intervenes when a candidate action would violate them. Because this mechanism is external to the base policy, shielding can be combined naturally with DRL.

\subsection{Single-agent shielding}

In the single-agent setting, shielding can be preemptive or post-posed. Preemptive shielding restricts the available action set before the agent acts, whereas post-posed shielding replaces an unsafe chosen action only when necessary \cite{elsayedaly2021safemultiagentreinforcementlearning}. In both cases, the shield is an external logic component that enforces safety with minimal interference.

\subsection{Multi-agent shielding}

In multi-agent settings, unsafe outcomes arise from joint actions. ElSayed-Aly et al. distinguish centralized and factored multi-agent shielding \cite{elsayedaly2021safemultiagentreinforcementlearning}. Centralized shields monitor the full joint state and joint action, but suffer from exponential scaling in the number of agents. Factored shielding mitigates this by decomposing the problem into smaller local (decentralized) shields.

\subsection{Shielding limitations}

Classical shielding is typically deterministic and rejection-based: unsafe actions are blocked at runtime. In realistic autonomous driving, this is restrictive \cite{yang2023safereinforcementlearningprobabilistic}. It assumes perfect state information and binary safety, even though real observations are noisy and partial.
\\
\\
Rejection-based shielding also creates a learning mismatch \cite{minihackexp}. If unsafe actions are silently replaced, the policy is updated as if its original actions had been executed, which biases learning and weakens convergence guarantees.

\section{Probabilistic Logic Shields (PLS)}

Probabilistic Logic Shields (PLS) \cite{jansen2019safereinforcementlearningprobabilistic} address these limitations by combining DRL policies with probabilistic logical reasoning. A neural policy proposes an action distribution, symbolic sensors extract probabilistic predicates from the same observation, and safety rules are encoded as probabilistic logic programs \cite{yang2023safereinforcementlearningprobabilistic}. Instead of rejecting actions outright, the shield reweights the policy distribution toward safer actions. This allows uncertainty, noisy sensing, and partial observability to be handled more naturally.
\\
\\
A key advantage is that PLS remains differentiable and compatible with gradient-based DRL. The effect of the shield can therefore enter the policy update itself rather than remaining only an external correction layer \cite{minihackexp}.

\subsection{Single-agent PLS}

Empirical studies in single-agent MiniHack environments show that PLS can reduce unsafe behavior during training, speed up convergence by reducing unproductive exploration, and improve generalization under environment changes \cite{minihackexp}. These results make PLS a promising candidate beyond the single-agent setting.

\subsection{Multi-agent PLS}

Although PLS was introduced in the single-agent setting, Chatterji and Acar extend the idea to multi-agent reinforcement learning through Shielded Multi-Agent Reinforcement Learning (SMARL) \cite{chatterji2025thinksmartactsmarl}. In SMARL, each agent combines its base policy with symbolic sensor information in a local probabilistic logic shield, yielding a reweighted safe policy. Because multiple agents interact, zero violations at every step are no longer guaranteed as in simpler single-agent settings \cite{minihackexp}. However, Chatterji and Acar show that the joint shielded policy is at least as safe as the unshielded joint policy, and strictly safer when any local policy violates its safety constraints \cite{chatterji2025thinksmartactsmarl}.
\\
\\
SMARL is evaluated on coordination games, social dilemmas, mixed-motive problems, and spatio-temporal cooperation tasks such as Stag Hunt, the Extended Public Goods Game, and Centipede \cite{chatterji2025thinksmartactsmarl}. The results suggest that multi-agent PLS can act as an equilibrium-selection mechanism, reduce violations, and strengthen cooperative behavior. The same work also studies asymmetric shielding and shows that shielding only a subset of agents can still induce safer behavior in the overall population \cite{chatterji2025thinksmartactsmarl}.

\section{Urban traffic simulators}

Simulation is central to autonomous driving research because it enables controlled experiments on perception, planning, and interaction. Popular simulators such as CARLA \cite{dosovitskiy2017carlaopenurbandriving}, SUMO \cite{8569938}, AirSim \cite{shah2017airsimhighfidelityvisualphysical}, and HighwayEnv \cite{highway-env} provide realistic traffic scenarios and physics, but they are less suited to flexible logical abstraction and explicit rule-based modeling.
\\
\\
Some more recent tools, such as SCENIC \cite{elmaaroufi2024scenicnlgeneratingprobabilisticscenario}, introduce richer formal specification layers, but safe RL and multi-agent shielding still benefit from environments in which rules and agent configurations can be manipulated directly and systematically.

\subsection{LogiCity simulator}

Li et al. introduce LogiCity \cite{li2025logicityadvancingneurosymbolicai}, a simulator designed around customizable first-order logic abstractions rather than fixed scripted traffic behavior. Agents and environment entities are represented through semantic and spatial predicates, and behavior is governed by logical rules. Because these abstractions are reusable across maps and scenarios, LogiCity is well suited to studying rule-based traffic interaction and neurosymbolic safety mechanisms.
\\
\\
LogiCity also uses rule-based agents for non-RL participants such as pedestrians, allowing them to adapt consistently to shared traffic constraints \cite{li2025logicityadvancingneurosymbolicai}. Together with its support for long-horizon and multi-agent traffic tasks, this makes it a suitable testbed for the thesis focus on probabilistic logic shielding in multi-agent autonomous driving.


% \chapter{Background}
% \label{cha:2}

% \section{Autonomous Driving}

% An autonomous driving (AD) system typically consists of the following layers: sensor, perception, scene representation, planning/deciding and control \cite{kiran2021deepreinforcementlearningautonomous}. It is a pipeline that converts raw sensor data into control actions. The sensor layer integrates sensors such as cameras, radar, LiDAR, etc. to capture raw data of the environment. The perception layer constructs an intermediate representation of the scene using the raw data from the sensor layer, capturing information such as road structure, traffic participants, and signals. The scene representation layer abstracts it further trough sensor-fusion methods, a current scene representation of the local environment of the vehicle is constructed and detected objects and other vehicles are mapped on that representation. Based on this contextual information, the deciding/planning layer generates feasible trajectories respecting vehicle dynamics and constraints. Finally, the controlling layer takes the necessary actions such as steering, acceleration, and braking, to realize and follow the decisions and planning made in the previous layer.
% \\
% \\
% The decision-making component of AD systems is a complex challenge by itself. Most of the decision-making problems cannot be solved by only supervised learning. This is because an autonomous vehicle’s actions influence future states and involve uncertainty and sequential dynamics \cite{kiran2021deepreinforcementlearningautonomous}. Thus, driving decisions are modeled as a sequential decision process, which carries it to the field of reinforcement learning (RL). The agent learns an optimal policy through interactions with the environment. Of course, applying RL to autonomous vehicles has become a popular problem as vehicles must learn to make optimal decisions in combinatorial large spaces (traffic scenarios, road conditions, etc.) that are difficult to program explicitly.
% \\
% \\
% A key challenge in autonomous driving is the dynamic and interactive nature of traffic. Besides the autonomous vehicle itself, there are other agents present (other vehicles, pedestrians, etc.). In other words, the autonomous vehicle operates in an environment shared with other agents, whose behavior affects the optimal actions of the ego vehicle. Each vehicle’s state and safety depend not only on its own actions, but also on the actions of the others, making the environment highly non-stationary and unpredictable \cite{zhang2024multiagentreinforcementlearningautonomous}. Think of an intersection scenario or during a lane merge, where an autonomous car must anticipate and react to the decisions of human-driven or other autonomous cars. Thus, unlike simple obstacle avoidance or geometric path planning, driving in traffic is fundamentally a multi-agent, strategic task, which elevates the problem to multi-agent reinforcement learning (MARL). 

% \section{Multi-agent autonomous driving}

% Real-world urban traffic is a multi-agent case: autonomous vehicles (AVs) interact with each other, pedestrians, cyclists etc. According to the survey on Multi-Agent Reinforcement Learning (MARL) for autonomous driving \cite{zhang2024multiagentreinforcementlearningautonomous}, traditional single-agent RL approaches can produce more self-centered or aggressive behaviors. This is because they optimize only the reward of the agent itself. Multi-agent systems (MAS) however, have to explicitly model interactions, mutual influence, and coordination requirements between agents, which aligns more with real-world traffic systems. 
% \\
% \\
% Autonomous driving environments in MARL are typically formalized as Decentralized Partially Observable MDPs (Dec-POMDPs). In these environments, each agent observes only a limited field of view and interacts with uncertainty, which aligns closely with the reality. And depending on the task, the agents can have cooperative, competitive, or mixed-motivation (balance between the two) interactions. The survey \cite{zhang2024multiagentreinforcementlearningautonomous} also states the difference between different training paradigms, like Centralized Training with Decentralized Execution (CTDE) and fully Decentralized Training and Execution (DTDE). The CTDE approach allows a centralized critic to exploit global information during training and preserves decentralized operation at runtime. DTDE on the other hand, removes the central coordination entirely and is more practical for large-scale traffic systems.

% \section{Deep Reinforcement Learning}

% Autonomous driving is a sequential decision-making problem where the vehicle has to choose actions that influence future observations and interactions with other agents in a dynamic traffic environment. These type of problems are typically modeled using Markov Decision Processes (MDPs). The process is an agent–environment interaction loop where the vehicle is the agent that constantly observes the state of its surroundings (the environment) and selects control actions using a certain policy. Deep reinforcement learning (DRL) is a promising approach for handling decision-making tasks in autonomous driving, as DRL algorithms combine reinforcement learning with deep neural networks. Using deep neural networks enables agents to learn complex policies from high-dimensional inputs efficiently.
% \\
% \\
% % Abbrevation of the methods
% DRL can be applied to tasks such as lane keeping, lane changing, merging, intersection handling etc. In both discrete and continuous settings, methods like DQN variants, actor–critic approaches (e.g., DDPG, A3C, PPO), and hybrid DRL-plus-control architectures have been used \cite{kiran2021deepreinforcementlearningautonomous}.
% \\
% \\
% However safety concerns arise on safety critical systems with DRL, because in those systems, even a single failure cannot be tolerated and learning must be assured to always satisfy the safety constraints \cite{isele2019safereinforcementlearningautonomous}. Unsafe exploration, rare edge cases etc. can hereby lead to problematic consequences. Therefore, DRL approaches are not enough on their own and must be improved with explicit safety constraints, domain knowledge, and formal guarantees. This motivates the exploration of safe reinforcement learning techniques.

% \section{Safe Reinforcement Learning}
% A core difficulty in applying RL to autonomous driving (and other safety-critical domains) is the lack of safety guarantees. Standard RL agents learn through trial-and-error. This means that they will inevitably experience (and cause) collisions or other unsafe events. The RL agent does not consider the safety of its taken actions, as long as the (final) reward is maximized. In the real world, such unsafe exploration is unacceptable. Safe reinforcement learning (Safe RL) aims to learn optimal policies while respecting safety constraints. In other words, the agent should avoid dangerous actions and states during the training and the deployment of the agent.

% \subsection{Formalization of safety in reinforcement learning}
% In Safe RL, \emph{constrained Markov decision process} (CMDP) is a widely-used formal framework. The goal of the agent is to maximize the expected return, while ensuring that one or more cost functions remain below the predefined thresholds \cite{gu2024reviewsafereinforcementlearning}. The cost functions represent the safety violations. In this way, the safety-performance trade-offs can be expressed. 
% \\
% \\
% In autonomous driving however, CMDP-based formulations face important limitations. Many safety requirements are state-dependant and instant decisions (e.g., collisions must never occur). CMDPs however typically enforce constraints in expectation over trajectories. Moreover, autonomous driving is a case of partial observability and involves non-stationary multi-agent interactions, making the reliable estimation and enforcement of safety costs very difficult. As a result, CMDP-based optimization by itself rarely provides strong, interpretable safety guarantees in realistic driving scenarios \cite{zhang2024multiagentreinforcementlearningautonomous}.

% \subsection{Categories of Safe RL methods}
% Safe RL methods are typically formalized using CMDPs as mentioned earlier, and they can be categorized based on how the safety is enforced. Summarizing the methods briefly based on the review on safe RL by Gu et al, (2024) \cite{gu2024reviewsafereinforcementlearning}:

% \paragraph{Reward-based and penalty-based methods.}
% Assure safety by adding penalty terms directly to the reward function, discouraging unsafe behavior during learning. While it is simple and widely used, these approaches have weak guarantees and depend heavily on the design of the rewards. In autonomous driving, designing the reward functions for DRL agents is still an open question  \cite{kiran2021deepreinforcementlearningautonomous}.

% \paragraph{Constrained optimization methods.}
% Instead of adding the constraints directly, the safety constraints are incorporated explicitly, which is commonly done via primal-dual or Lagrangian techniques. These methods have a stronger theoretical base, but the safety is usually only enforced in expectation. This still allows violations during exploration. 

% \paragraph{Control-theoretic methods.}
% These methods use tools such as Lyapunov functions, model predictive control, or control barrier functions to restrict the agent’s actions to a safe set of actions. Control-theoretic approaches can have strong guarantees, but require often accurate dynamics and an environment model. It is computationally expensive in multi-agent traffic.

% \paragraph{Formal methods–based methods.}
% These methods ensure safety by using mathematical logic and verification techniques. This helps to prove the guarantee that unsafe states cannot be reached. Because the safety is not learned, but written as formal rules. Formal methods struggle with continuous state spaces, high-dimensional dynamics and multi-agent interactions. This scalability problem of course directly applies to multi-agent autonomous driving environments.

% \paragraph{Gaussian Process–based methods.}
% These methods address safety by modeling the uncertainty in dynamics or cost functions. The core idea is that the agent behaves conservatively and avoids that region when the agent is uncertain about the consequences of an action in that region. Also these methods suffer from the scalability problem.

% \paragraph{Safety-layer and intervention-based methods.}
% These methods enforce safety at runtime by filtering or correcting actions proposed by the learned policy. The safety enforcement is separated from the learning component, which provides strong safety guarantees. A popular example of these methods are the shielding-based approaches. 

% \section{Shielding}

% Shielding is a Safe RL method \cite{alshiekh2017safereinforcementlearningshielding}, which is part of the Safety-layer and intervention-based methods. Shielding adds a formally synthesized runtime safety layer around the learning agent. The key idea is to encode safety requirements such as collision avoidance, lane-keeping, or respecting to traffic rules as temporal logic specifications. Then, the logic (shield) component monitors the agent's possible actions and intervenes if the action will violate the safety rules implemented in the shield. This mechanism can be combined with deep reinforcement learning without modifying the internal RL loop structure, making shielding a suitable candidate for autonomous driving.

% \subsection{Single-agent shielding}
% In the single-agent setting, the shield can be integrated in two ways. In preemptive shielding, the shield filters the available action set before the agent chooses an action. Therefore, only safe actions are exposed to the agent. In post-posed shielding, the shield checks the action proposed by the learning agent and replaces it only if it is unsafe \cite{elsayedaly2021safemultiagentreinforcementlearning}. in both cases, the RL process does not change internally, as mentioned earlier. The shield is an external logic component that enforces minimum interference, meaning the shield only intervenes when necessary (e.g., violation of safety rules). This allows for autonomous driving to incorporate safety rules (e.g., “no collision,”, “stay within drivable lanes”) easily to the DRL training process of the autonomous vehicle. 

% \subsection{Multi-agent shielding}

% In multi-agent settings, unsafe outcomes arise from the joint actions of multiple agents. ElSayed-Aly et al. extend shielding to multi-agent reinforcement learning (MARL) and have mentioned two ways of multi-agent shielding \cite{elsayedaly2021safemultiagentreinforcementlearning}. First of all, centralized shields. Centralized shields observe the joint state and the joint action of all agents at each time step, blocking or modifying unsafe joint maneuvers. A disadvantage of this method is that for many agents, number of joint states and joint actions increases exponentionally, which makes the computational cost of synthesizing these shields expensive, a classic scalability problem. The second method, which solves this scalability problem, is factored shielding. This method decomposes the joint state space into local regions (e.g., road segments or intersections) and synthesizes multiple smaller shields. Each of those shields is responsible for another local region, and agents can join or leave those local shields while moving trough the environment. 

% \subsection{Shielding limitations}

% The shielding methods mentioned above are deterministic rejection-based shielding methods, as it blocks unsafe actions at runtime. Those shields suffer from important limitations in realistic autonomous driving settings \cite{yang2023safereinforcementlearningprobabilistic}. First of all, perfect state information and binary safety (action is safe or unsafe) is assumed. This is unrealistic due to sensor data being noisy and having partial observability in reality. As a result, rejection-based shields cannot exploit uncertainty in information and therefore cannot distinguish between actions that are certainly unsafe and those that are as good as safe.
% \\
% \\
% Second, rejection-based shielding introduces a mismatch between the learning agent and the safety mechanism \cite{minihackexp}. Unsafe actions proposed by the policy are silently rejected by the shield and replaced by safe alternatives. However at the same time, the agent continues to update its policy as if actions were sampled directly from its own policy distribution. The agent receives biased feedback, because the agent is unaware of which actions were rejected. As a result of this mismatch, the convergence guarantee to an optimal policy in the RL process does not hold anymore. 

% \section{Probabilistic Logic Shields (PLS)}

% Probabilistic Logic Shields (PLS) \cite{jansen2019safereinforcementlearningprobabilistic} extend shielding by combining deep reinforcement learning policies with probabilistic logical reasoning. This allows the agent to enforce safety constraints and handles uncertainty as well. In PLS, the base policy is produced by a neural network as in standard DRL. Then, the same observation is passed through a set of symbolic sensors that extract probabilistic predicates about the environment like the likelihood of an obstacle being present close to the vehicle. The safety rules are defined as probabilistic logic programs \cite{yang2023safereinforcementlearningprobabilistic}, where unsafe situations are defined as rules in function of state predicates and the candidate actions. Instead of rejecting the actions, the shield renormalises the policy distribution by reweighting the policy probabilities in favor of the safe actions. By reasoning over probability distributions instead of making binary assumptions, PLS can handle noisy sensor data, uncertainty, and is easy to use with gradient DRL algorithms.
% \\
% \\
% A key advantage of PLS is that shielding becomes differentiable and compatible with the gradient-based DRL algorithms. With PLS, the intervention of the shield can also be included in the gradient updates. This way, the DRL agent sees the changes made by the shield. The PLS component can be seen as the continuous correction layer \cite{minihackexp}. 

% \subsection{Single-agent PLS}
% Empirical evaluations on MiniHack (single-agent) environments \cite{minihackexp} has shown several benefits of the PLS method compared with DRL methods using no shield at all. First of all, PLS agents avoid unsafe states during the training. Second, The agent converges in much fewer steps. The reason of this is that the Shield component also eliminates unproductive exploration. Finally, the agent reaches better generalization under environment changes. These results makes PLS a promising method to also extend it to multi-agent systems. 

% \subsection{Multi-agent PLS}
% While PLS was originally developed for single agent Safe RL, Chatterji and Acar have introduced Shielded Multi-Agent Reinforcement Learning (SMARL) \cite{chatterji2025thinksmartactsmarl}. SMARL is a general framework in which each agent interacts with its own probabilistic logic shield during training and deployment. In SMARL, each agent $p_i$ computes a base policy $\pi_i(s_i^t)$ and forwards this policy. The shield combines this base policy together with the sensor data to a ProbLog program $T_i$. The shield then produces a reweighted safe policy $\pi_i^{+}(s_i^t)$ based on its safety rules. SMARL does not guarantee does not guarantee zero violations at every time step, as the minihack experiments in single-agents had shown \cite{minihackexp}, because of the fact that multiple agents are involved. However Chatterji and Acar have proven that the joint shielded policy $\vec{\pi}^{+}$ is at least as safe as the unshielded joint policy $\vec{\pi}$, and it is strictly safer whenever any agent’s base policy violates its local safety constraints \cite{chatterji2025thinksmartactsmarl}.
% \\
% \\
% SMARL is tested on multi-agent simple coordination games, social dilemmas, mixed-motive and spatio-temporal cooperation tasks (e.g., Stag Hunt, Extended Public Goods Game, Centipede) \cite{chatterji2025thinksmartactsmarl}. The results have shown that first of all, PLS in multi-agent systems can function as an equilibrium-selection mechanism. Second, shielded agents performed fewer constraint violations and stronger cooperative behavior than unshielded baselines. Also, asymmetric shielding is analyzed whereby not all the agents get a shield, but a subset of them \cite{chatterji2025thinksmartactsmarl}. The experiments have shown that even when only a subset of agents is shielded, their behavior can stimulate unshielded agents toward behaving safer compared to using no shield for none of the agents. 

% \section{Urban traffic simulators}

% Simulation environments play a central role in the study of autonomous driving. Simulators allow controlled experimentation with perception, planning, and interaction dynamics. Popular AD simulators such as CARLA \cite{dosovitskiy2017carlaopenurbandriving}, SUMO \cite{8569938}, AirSim \cite{shah2017airsimhighfidelityvisualphysical} and HighwayEnv \cite{highway-env} provide decent rendering, realistic physics and configurable traffic scenarios. However, these environments typically rely on fixed rule sets and scenario-specific dynamics. Thus, the possibilities with those simulators are limited for explicit abstraction or flexible configuration of agents and rules. 
% \\
% \\
% Some more recent tools introduce formal specification layers such as SCENIC \cite{elmaaroufi2024scenicnlgeneratingprobabilisticscenario}, which enables scenario generation using temporal-logic constraints. However also this simulator is situating near predefined traffic elements and task-specific safety conditions. Research on safe reinforcement learning and multi-agent shielding requires environments where logical rules, abstractions, and agent configurations can be worked with flexibly. This allows safety mechanisms to be evaluated not only as constraint checkers but also as part of the learning and coordination process itself.

% \subsection{LogiCity simulator}

% Recently, Li et al. introduced the LogiCity simulator \cite{li2025logicityadvancingneurosymbolicai}. This simulator directly adresses those limitations by modeling urban environments through customizable first-order logic (FOL) abstractions rather than fixed scripted behaviors. The agents and evironment entities are represented via semantic and spatial predicates like \emph{IsAmbulance(X)} and \emph{IsClose(X,Y)}. And to lead the agent's behavior, those predicates are combined with rules. The abstractions are independent of specific agents or maps, which allows generalization and creation of different urban scenarios by simply reusing the same logical abstractions and rules. 
% \\
% \\
% For simulating the non-RL agents like pedestrians, LogiCity uses rule-based agents instead of manually scripted behavior trees. This allows the agents to automatically adapt actions to maintain consistency with the global shared constraints like yielding to emergency vehicles for example. Furthermore, the simulator supports long-horizon sequential decision-making tasks and multi-agent interaction scenarios as well.



\chapter{Methodology}

\section{Problem Formulation}
This thesis investigates probabilistic logic shielding (PLS) in multi-agent environments for safe autonomous driving, with a particular focus on \emph{safe task completion} under shared traffic rules. The central objective is not only to test whether PLS reduces unsafe behavior, but to study how shielding changes the balance between successful completion, direct failure, and timeout as multiple RL-controlled cars interact in the same traffic scene. In particular, the thesis examines how PLS behaves as multi-agent interaction pressure increases, how decentralized shield styles trade off conservatism against efficiency, whether centralized shielding can resolve coordination deadlocks that remain difficult under decentralized local shielding, and whether homogeneous multi-agent traffic control is better served by shared-policy or separate-policy training.

\subsection{Safe Multi-Agent Navigation in Urban Traffic Scenarios}
The considered task is a navigation problem in a stylized urban traffic environment implemented in LogiCity, with a particular focus on intersection scenarios. Intersections are especially suitable for this thesis because they concentrate the main coordination challenges of urban driving in a compact setting: vehicles must account for right-of-way, yielding, shared occupancy, and pedestrian interaction rather than following their routes independently.

\begin{figure}[H]
    \centering
    \includegraphics[width=0.50\textwidth]{figures/exp1_k2_val_ep0_frames/step_4.png}
    \caption{Example intersection scenario used in the thesis.}
    \label{fig:exp1_intersection_scenario}
\end{figure}

\noindent The experiments therefore use compact intersection-centered traffic scenes. Figure~\ref{fig:exp1_intersection_scenario} shows a representative example in which one or more controllable private cars act as the RL agents, while pedestrians serve as background agents. Start and goal positions are assigned so that the central intersection must be traversed, forcing interaction with both other cars and pedestrians.
\\
\\
The Safe Path Following (SPF) task is defined as follows \cite{li2025logicityadvancingneurosymbolicai}: all controllable cars must reach their respective goal positions safely and successfully while avoiding collisions and obeying the specified traffic rules. In particular, cars must not collide with pedestrians or other controllable cars, and they must not violate rules related to yielding, right-of-way, or unsafe intersection occupancy. The task is therefore not merely path planning, but multi-agent traffic coordination under safety constraints.

\subsection{State, Action, and Observation Spaces}
The joint traffic environment is defined over the full simulator state, which includes the map layout, road and sidewalk structure, intersection regions, all agent positions, planned routes, goals, and traffic-relevant logical relations. This full state is used internally by the simulator and by the rule machinery, but it is not given directly to the reinforcement learning policy. Instead, each controllable car acts from a local symbolic observation centered on itself.

\begin{table}[H]
\centering
\small
\caption{State, action, and observation design of the task.}
\label{tab:state-action-observation}
\begin{tabularx}{\linewidth}{@{}p{0.2\linewidth} X@{}}
\toprule
\textbf{Component} & \textbf{Description} \\
\midrule
Environment state & Full simulator state, including map structure, agent positions, routes, goals, and logical traffic relations. Available to the simulator and rule system, but not directly to the learned policy. \\
\addlinespace
Action space & Discrete four-action speed control: \texttt{Slow}, \texttt{Normal}, \texttt{Fast}, and \texttt{Stop}. \\
\addlinespace
Observation space & Flattened grounded logical vector over a local field of view. Unary predicates contribute one value per visible entity; binary predicates contribute one value per ordered entity pair. \\
\addlinespace
Observation scope & One full intersection region, with at most three visible entities. \\
\bottomrule
\end{tabularx}
\end{table}

The RL policy acts in a discrete four-action space:
\[
\mathcal{A} = \{\texttt{Slow}, \texttt{Normal}, \texttt{Fast}, \texttt{Stop}\}.
\]

\begin{table}[H]
\centering
\small
\caption{Discrete action space of each controllable car.}
\label{tab:exp1-action-space}
\begin{tabularx}{\linewidth}{@{}l X@{}}
\toprule
\textbf{Action} & \textbf{Effect} \\
\midrule
\texttt{Slow} & One-cell movement in one of the four cardinal directions. \\
\texttt{Normal} & Two-cell movement in one of the four cardinal directions. \\
\texttt{Fast} & Three-cell movement in one of the four cardinal directions. \\
\texttt{Stop} & No movement. \\
\bottomrule
\end{tabularx}
\end{table}

\noindent Thus, the RL problem is fundamentally a speed-mode decision problem under logical safety constraints.
\\
\\
We do not use visual observations. Each RL-controlled agent receives a flattened grounded logical state vector over its local field of view.
\\
\\
\noindent The grounded predicates include semantic type predicates, traffic-state predicates, and relational predicates, such as whether an entity is a car or pedestrian, whether it is at or in an intersection, whether it has higher priority than the ego vehicle, and whether it is collision-close. The full predicate list is provided in Appendix~\ref{app:predicate-definitions}.
\\
\\
Figure~\ref{fig:exp1_intersection_scenario} gives a simple example. If the controllable car approaching from the south is treated as the ego agent, then the two pedestrians in the central crossing area are part of its local observation, whereas the car approaching from the north is still outside its field of view. The observation is therefore intentionally local: it captures only the immediate interaction structure around the ego agent.
\\
\\
For this case, a schematic local entity set is
\[
[\texttt{ego}, \texttt{Pedestrian\_1}, \texttt{Pedestrian\_2}].
\]
\noindent The approaching northern car is not yet included because it is still outside the ego field of view.

\begin{table}[H]
\centering
\small
\caption{Example grounded observation entries for the southern ego car.}
\label{tab:exp1-grounded-observation-example}
\begin{tabular}{@{}ll@{}}
\toprule
\textbf{Ground atom} & \textbf{Value} \\
\midrule
\texttt{IsCar(ego)} & 1 \\
\texttt{IsAtInter(ego)} & 1 \\
\texttt{IsInInter(ego)} & 0 \\
\texttt{IsPedestrian(Pedestrian\_1)} & 1 \\
\texttt{IsInInter(Pedestrian\_1)} & 1 \\
\texttt{IsPedestrian(Pedestrian\_2)} & 1 \\
\texttt{IsInInter(Pedestrian\_2)} & 1 \\
\texttt{HigherPri(Pedestrian\_1, ego)} & 0 \\
\texttt{HigherPri(Pedestrian\_2, ego)} & 0 \\
\texttt{CollidingClose(Pedestrian\_1, ego)} & 0 \\
\texttt{CollidingClose(Pedestrian\_2, ego)} & 0 \\
\bottomrule
\end{tabular}
\end{table}

\noindent The policy does not receive these atoms symbolically. Instead, they are flattened into a fixed-order numeric vector, for example
\[
[1,1,0,\;1,1,\;1,1,\;0,0,\;0,0,\dots],
\]
\noindent where the exact ordering is determined by the implementation. The input is therefore numeric, but each dimension still retains a clear logical interpretation.

\subsection{Observation Noise Model}
The thesis also considers an uncertain symbolic observation regime. This noise is part of the environment-side observation model rather than of PLS itself: the LogiCity dynamics remain unchanged, but the symbolic information provided to the policy and shield is corrupted to mimic imperfect sensing.
\\
\\
The implemented noise process is structured rather than i.i.d. It distinguishes between type and relation predicates, treats ego-related facts differently from non-ego facts, and increases corruption for crowded or occluded local observations. As a result, relational traffic predicates such as collision proximity, intersection occupancy, and priority can become more uncertain than simple type labels such as whether an entity is a car or a pedestrian.
\\
\\
This matters because the most safety-relevant predicates are relational. The shield does not define this uncertainty model; it simply consumes the noisy grounded observation provided by the environment. The full specification is given in Appendix~\ref{app:observation-noise-model}.

\subsection{Agents}
The Safe Path Following task contains two classes of agents: RL agents, whose behavior is learned and evaluated in this thesis, and background agents, which remain controlled by LogiCity's rule-based expert planner.
\\
\\
Background agents are not replayed along fixed trajectories. Instead, the expert planner repeatedly evaluates the current local traffic configuration and selects admissible route-following actions from logical traffic predicates \cite{li2025logicityadvancingneurosymbolicai}. They therefore react online to conditions such as intersection occupancy, relative priority, and collision-related relations.
\\
\\
In all experiments, pedestrians are always background agents, while all cars are RL agents. The learning problem is therefore centered on the interacting car population, while pedestrians provide structured external traffic interaction.
\\
\\
For simplicity, the thesis uses only two traffic-participant types: normal cars and pedestrians. Although LogiCity supports richer semantic roles such as ambulances, police cars, and buses \cite{li2025logicityadvancingneurosymbolicai}, these are not used here because the focus is on multi-agent autonomous navigation rather than semantic traffic heterogeneity.

\subsection{Safety Objective and Efficiency Objective}
The task combines a safety objective with an efficiency objective. Safety means avoiding both terminal and non-terminal unsafe behavior. Terminal unsafe behavior corresponds to realized collisions, which end the episode immediately, while non-terminal violations include near-collision cues, right-of-way violations, and simultaneous intersection occupancy.
\\
\\
Efficiency means reaching the assigned goal positions with as little unnecessary delay as possible. In the current reward design, this is encouraged through a positive terminal goal reward and a progress-based shaping reward rather than through an explicit per-step time penalty.
\\
\\
This distinction is important for interpreting the role of the shield. A conservative controller may reduce unsafe behavior but still perform poorly through repeated yielding and timeout, whereas a more aggressive controller may improve progress at the cost of more unsafe interaction. The methodology therefore treats safe completion as the desired outcome and views collisions, traffic-rule violations, and timeouts as distinct manifestations of failure or inefficiency.

\iffalse
\subsection{Difficulty Scaling Through the Number of Agents \texorpdfstring{$K$}{K}}
To study how shielding behaves under increasing interaction complexity, the benchmark varies the number of controllable cars through the parameter $K$. In this thesis, $K$ denotes the number of private car agents that are simultaneously controlled by the shared reinforcement learning policy. The compact Experiment~1 benchmark considers three primary settings: $K=1$, $K=2$, and $K=3$. Across these settings, the map and the general traffic structure remain fixed, while the number of interacting controllable vehicles increases.
\\
\\
Importantly, the increase in $K$ does not simply add more copies of the same isolated task. Each additional car introduces more possible right-of-way relations, more potential occupation conflicts at the central intersection, and more opportunities for conservative waiting behavior to propagate through the system. The pedestrian components remain present as structured sources of crossing interaction, while the vehicle roster grows from one controllable car in the $K=1$ setting, to two controllable cars in $K=2$, and three controllable cars in $K=3$. This design makes $K$ a direct measure of multi-agent interaction density within the same compact scenario family.
\\
\\
The use of $K$ as the main difficulty axis is central to the experimental question of the thesis. If probabilistic shielding is too conservative, one expects increasing $K$ to produce more yielding, more stopping, and eventually more timeout-dominated behavior. If the shield is too permissive, increasing $K$ may instead expose more unsafe interactions. By scaling $K$ while keeping the general scenario structure fixed, the benchmark isolates how shield style affects the balance between safety preservation and efficient coordination under growing multi-agent complexity.
\fi

\section{Task Rules}
\subsection{Task Rules as Logical Constraints}
The task rules are defined as explicit logical formulas over grounded traffic predicates. In the current thesis setup, the rule system evaluates predicates such as \texttt{IsAtInter}, \texttt{IsInInter}, \texttt{HigherPri}, \texttt{CollidingClose} and \texttt{Collision}, and then determines whether the current state-action combination satisfies or violates the task requirements.
\\
\\
The central modeling idea is stop-style traffic logic: whenever a rule-relevant hazard condition holds, the safe traffic action is \texttt{Stop(entity)}. Such hazard conditions include pedestrian-crossing conflicts in the same intersection, entering an intersection while another agent is already occupying it, approaching the same intersection together with another car that has higher priority, simultaneous same-intersection occupancy by two cars, and imminent collision-close situations. The task rules therefore do not describe safety in terms of geometric free space alone, but in terms of logical traffic relations that encode right-of-way, shared intersection occupancy, and conflict proximity.

\begin{table}[H]
\centering
\small
\caption{Core task-rule predicates in the SPF task.}
\label{tab:task-rule-predicates}
\begin{tabularx}{\linewidth}{@{}l X@{}}
\toprule
\textbf{Predicate} & \textbf{Role in the task rules} \\
\midrule
\texttt{IsAtInter} & Indicates that an agent is at the boundary of an intersection and is about to enter or yield at that region. \\
\texttt{IsInInter} & Indicates that an agent is currently occupying the interior of an intersection region. \\
% \texttt{SameInter} & Indicates that two agents are associated with the same intersection region, so that intersection-conflict rules apply between them. \\
\texttt{HigherPri} & Encodes right-of-way asymmetry between two traffic participants. \\
\texttt{CollidingClose} & Encodes an imminent collision cue that requires immediate stopping but does not itself represent a realized collision. \\
\texttt{Collision} & Encodes a realized collision state and is used for terminal failure. \\
\texttt{IsParked} & Marks inactive parked cars so that they are excluded from active conflict and failure evaluation. \\
\bottomrule
\end{tabularx}
\end{table}

\noindent Parked agents are excluded from rule triggering. The current implementation treats cars that are still parked at spawn or already parked after reaching their goal as inactive for both violation counting and terminal collision evaluation. This prevents task failures from being generated by inactive parked cars and keeps the task semantics focused on active traffic interaction.

\subsection{Non-Terminating Traffic Rules}
Most traffic-rule violations are modeled as non-terminating task rules. Violating such a rule produces a negative task signal but does not immediately end the episode. This design is important because it separates unsafe or undesirable traffic behavior from realized catastrophic failure.
\\
\\
\begin{table}[H]
\centering
\small
\caption{Non-terminating task rules.}
\label{tab:nonterminal-task-rules}
\begin{tabularx}{\linewidth}{@{}l X@{}}
\toprule
\textbf{Rule} & \textbf{Meaning} \\
\midrule
\textbf{Collision Close} & If an imminent collision cue is present, the ego agent is required to stop. Ignoring this cue produces a traffic-rule violation. \\
\textbf{Intersection Right-of-Way} & If the ego agent faces a pedestrian-crossing conflict, an occupied intersection conflict, or a higher-priority same-intersection conflict, it is required to stop. Failing to stop counts as a non-terminal traffic-rule violation. \\
\textbf{Simultaneous Intersection Entry} & If two cars occupy the same intersection at the same time, this is counted as a non-terminal rule violation even if no realized collision has yet occurred. \\
\bottomrule
\end{tabularx}
\end{table}

\noindent These non-terminating rules are deliberately traffic-semantic. They are meant to capture situations in which the ego car behaves unsafely or unlawfully, even when the episode has not yet reached an irreversible terminal state. This distinction is especially important for studying PLS, because a shield may reduce unsafe interaction patterns even before it prevents actual collisions.

\subsection{Terminating Failure Condition}
The only terminating task rule in the current setup is the realized \texttt{Collision} predicate. When the environment detects that an active agent is in a realized collision state with another active agent, the corresponding task rule is violated and the episode terminates with failure.
\\
\\
This means that the task distinguishes clearly between \emph{traffic-rule danger cues} and \emph{realized physical failure}. Near-collision conditions, right-of-way violations, and simultaneous intersection occupancy are treated as serious but non-terminal events, whereas a realized collision is treated as the catastrophic endpoint of unsafe interaction.

\begin{table}[H]
\centering
\small
\caption{Episode outcomes induced by the task rules.}
\label{tab:task-rule-outcomes}
\begin{tabularx}{\linewidth}{@{}l X@{}}
\toprule
\textbf{Outcome} & \textbf{Condition} \\
\midrule
\textbf{Success} & All controllable cars reach their assigned goals without triggering a terminal failure. \\
\textbf{Failure} & A realized collision occurs and activates the terminal \texttt{Collision} rule. \\
\textbf{Timeout} & The episode horizon is exceeded before all controllable cars complete their routes. \\
\bottomrule
\end{tabularx}
\end{table}

\section{Probabilistic Logic Shielding}
\subsection{Overview of the Shielding Framework}
The shielded method used in this thesis instantiates a hybrid, probabilistic-logic shield in LogiCity, inspired by Probabilistic Logic Shielding (PLS) \cite{yang2023safereinforcementlearningprobabilistic}. ProbLog provides a fixed probabilistic action-program interface, while a traffic template computes the domain-specific safety scores from grounded symbolic observations. PLPG is the training algorithm used to optimize the shielded policy. The implemented shield therefore consists of three components: a probabilistic policy representation, a state-dependent fact set, and a traffic-safety template that maps those facts to action-level safety estimates.
\\
\\
The first component is the policy representation \(\Pi_s\), encoded as an annotated disjunction inside the logic program. For a given state \(s\), the neural policy is represented as
\[
\Pi_s =
\left\{
\begin{array}{l}
\pi_\theta(\texttt{slow}\mid s)::\texttt{act(slow)};\\
\pi_\theta(\texttt{normal}\mid s)::\texttt{act(normal)};\\
\pi_\theta(\texttt{fast}\mid s)::\texttt{act(fast)};\\
\pi_\theta(\texttt{stop}\mid s)::\texttt{act(stop)}
\end{array}
\right\}.
\]
\noindent The second component is the probabilistic fact set \(H_s\), which abstracts the current traffic state into shield-relevant symbolic information. The third component is the safety template itself. In the current implementation, a minimal ProbLog program is compiled once and reused at runtime, while the detailed traffic-safety computation is carried out by the template from the grounded predicate values. Conceptually, this can still be summarized as one shield program per state \cite{yang2023safereinforcementlearningprobabilistic},
\[
\mathcal{T}(s)=\mathcal{BK}\cup H_s\cup \Pi_s,
\]
\noindent but the important practical point is that the current shield uses a hybrid realization. ProbLog is used to represent the action distribution and to provide the fixed queried action structure of the shield program. The traffic semantics themselves, however, are not obtained by fully re-grounding a large symbolic program at every step. Instead, the grounded observation is passed into the soft-traffic template, which computes the relevant hazard quantities directly from traffic predicates such as \texttt{IsAtInter}, \texttt{IsInInter}, \texttt{SameInter}, \texttt{HigherPri}, and \texttt{CollidingClose}. These hazard quantities are then converted into action-level safety scores for \texttt{slow}, \texttt{normal}, \texttt{fast}, and \texttt{stop}, and those scores are finally used for policy reweighting. In this sense, ProbLog provides the probabilistic action skeleton, while the template supplies the domain-specific traffic-safety computation.

\subsection{Logical Grounding from Observations}
The probabilistic fact set \(H_s\) is derived directly from the grounded observation vector described earlier in this chapter. In other words, the shield does not observe the full simulator state independently of the policy. Instead, it consumes the same local symbolic observation that is already available to the RL-controlled car, and converts that local observation into shield-relevant logical facts.
\\
\\
In the current thesis setup, the shield abstraction is built from grounded predicates such as whether the ego car is at an intersection, whether it is already in the intersection, whether another visible entity belongs to the same intersection, whether that entity has higher priority, whether it is a pedestrian, and whether it is collision-close. These are precisely the facts required by the current traffic shield template. Each grounded fact is represented as a value in \([0,1]\), so that the shield reasons directly over the confidence-valued symbolic observation produced by the environment.

\subsection{Probabilistic Action-Safety Estimates}
For each candidate action \(a\), the traffic template computes an action-level safety score from the ego car's local grounded observation \(o^i\). We denote this score by
\[
q(o^i,a)\in[0,1].
\]
The fixed ProbLog program supplies the action-query interface, but the implemented action-dependent score is produced by the soft-traffic template after it evaluates the grounded traffic cues. Thus, \(q(o^i,a)\) is the template's local estimate of the safety of action \(a\), rather than the result of re-grounding and solving the full traffic rules in ProbLog at every decision step.
\\
\\
Using these action-level safety scores, the shield also computes the safety of the base policy,
\[
P_{\pi_\theta}(\texttt{safe}\mid o^i)
=
\sum_{a\in\mathcal{A}} \pi_\theta(a\mid o^i)\, q(o^i,a).
\]
\noindent This quantity measures how safe the unmodified policy is under the current symbolic traffic facts and the current background knowledge. It is therefore a state-dependent safety score for the base distribution itself rather than only for any single selected action.
\\
\\
An important consequence of this formulation is that rule satisfaction is treated probabilistically rather than only as a hard Boolean filter. The resulting safety scores remain soft because the shield combines confidence-valued symbolic facts with a probabilistic action distribution. As a result, the shield does not merely classify actions as admissible or inadmissible, but continuously expresses how safe each action is under the current observed traffic context.

\subsection{Action Reweighting Based on Safety Estimates}
The shielded policy is constructed by reweighting the base policy with the action-level safety scores \cite{yang2023safereinforcementlearningprobabilistic}:
\[
\pi_\theta^+(a\mid o^i)
=
\frac{\pi_\theta(a\mid o^i)\, q(o^i,a)}
{\sum_{a'\in\mathcal{A}} \pi_\theta(a'\mid o^i)\, q(o^i,a')}
=
\frac{\pi_\theta(a\mid o^i)\, q(o^i,a)}
{P_{\pi_\theta}(\texttt{safe}\mid o^i)}.
\]
\noindent This shielded distribution is the policy that is actually used for acting. Actions that are judged safer under the current symbolic traffic state receive relatively larger probability mass, while actions that are judged less safe are down-weighted.
\\
\\
Figure~\ref{fig:pls-worked-example} illustrates this reweighting with a concrete local traffic situation. The base policy initially prefers \texttt{fast}, but its low safety score reduces its probability after reweighting; the released probability mass is redistributed among the relatively safer actions.

\begin{figure}[H]
    \centering
    \includegraphics[width=0.90\textwidth]{figures/working_example PLS.png}
    \caption{Worked example of probabilistic action reweighting by the decentralized shield. The entries in each vector correspond to \texttt{slow}, \texttt{normal}, \texttt{fast}, and \texttt{stop}, respectively.}
    \label{fig:pls-worked-example}
\end{figure}

\noindent The safety of the shielded policy is \cite{yang2023safereinforcementlearningprobabilistic}
\[
P_{\pi_\theta^+}(\texttt{safe}\mid o^i)
=
\sum_{a\in\mathcal{A}} \pi_\theta^+(a\mid o^i)\, q(o^i,a).
\]
\noindent Thus, the shield acts as a probabilistic action-reweighting mechanism rather than as a purely external hard action blocker. This is important in the present setting because the method is intended not only to avoid obviously unsafe actions, but also to shape behavior continuously according to how safe different actions are under the current traffic context.

\subsection{PLPG Training Objective}
The previous subsection described how the shield modifies action selection at decision time. In the present implementation, however, the shield is also part of policy learning through the PLPG objective \cite{yang2023safereinforcementlearningprobabilistic}.
\\
\\
Rollout collection uses the shielded policy \(\pi_\theta^+\) rather than the raw base policy \(\pi_\theta\), so learning data is generated under shielded behavior. The PPO-style update is then also defined with respect to shielded probabilities, meaning that optimization is carried out directly in the reweighted policy space.
\\
\\
In addition, the learning objective includes an explicit safety regularizer \cite{yang2023safereinforcementlearningprobabilistic}, whose exact form is given in Appendix~\ref{app:safety-regularizer}. This encourages safer action distributions under the current symbolic traffic facts and makes the shield part of the learning computation rather than a post-processing layer applied only at execution time.
\\
\\
A key advantage of PLS is therefore that shielding remains compatible with gradient-based reinforcement learning: the differentiable reweighting performed by the shield is reflected directly in the policy updates \cite{yang2023safereinforcementlearningprobabilistic,minihackexp}.

\subsection{Traffic Shield Template}
The current experiments use one fixed \textbf{traffic shield template}. The semantic shield family is therefore held constant across all shielded runs; what changes later is only the strength of the safety weighting.
\\
\\
The ProbLog component is deliberately small and fixed. It is compiled once and supplies the candidate-action interface used by the shield:
\begin{lstlisting}[style=mystyle,language=Prolog,caption={Fixed ProbLog interface of the implemented shield.},label={lst:traffic-shield-skeleton}]
% Compiled once; policy probabilities replace 0.25 at runtime.
0.25::act(slow); 0.25::act(normal); 0.25::act(fast); 0.25::act(stop).

% Fixed safety-query interface.
safe.
query(safe).
query(act(slow)).
query(act(normal)).
query(act(fast)).
query(act(stop)).
\end{lstlisting}
\noindent The displayed action probabilities are compile-time placeholders. At runtime, they are replaced by the current PPO probabilities. This listing is an interface rather than the traffic-rule engine: the unconditional \texttt{safe} atom does not itself derive traffic safety. The actual action-dependent safety scores are computed by the soft-traffic template.
\\
\\
The template reads the local grounded observation and extracts intersection position, shared-intersection membership, relative priority, pedestrian presence, and collision-close cues. It combines these confidence-valued predicates into six continuous hazards: collision-close, pedestrian conflict, occupied intersection, higher-priority conflict, simultaneous intersection occupancy, and transition conflict. The largest relevant hazard and its derived traffic regime determine how strongly \texttt{slow}, \texttt{normal}, and \texttt{fast} are penalized; \texttt{stop} remains the safer option in critical traffic states. Appendix~\ref{app:soft-traffic-hazard-calculation} gives the complete implemented hazard and action-score calculation. This rule-inspired calculation is implemented procedurally rather than as a second ProbLog program.
\\
\\
The complete decentralized shield flow is shown in Figure~\ref{fig:traffic-shield-flow}. Each controlled car runs this procedure separately: it starts from its own local grounded observation, derives local hazard values, converts them into action-level scores \(q(o^i,a)\), and reweights its own PPO distribution. Consequently, the shield can react to locally visible conflicts but does not jointly select the actions of multiple RL-controlled cars.

\begin{figure}[H]
    \centering
    \includegraphics[width=\textwidth]{figures/Decentralized PLS.png}
    \caption{Pipeline of the implemented decentralized hybrid PLS shield for one controlled car \(i\). The same local procedure is applied independently to each RL-controlled car.}
    \label{fig:traffic-shield-flow}
\end{figure}

\noindent The resulting safety scores are action-dependent rather than purely binary. In benign states, movement actions remain relatively safe; under increasing interaction pressure, faster intersection entry is penalized first; and in critical states, \texttt{stop} is strongly preferred. Thus, the shield is probabilistic because it reweights a probability distribution using confidence-valued traffic information, and logic-grounded because its hazards are defined over symbolic traffic predicates. ProbLog provides the fixed action-program interface, while the traffic template provides the domain-specific graded safety computation used for reweighting.

\subsection{Conservative and Aggressive Shield Variants}
The conservative and aggressive shield variants differ by safety weighting, not by logical rule content. Both variants use the same observation grounding pipeline, the same traffic-hazard structure, and the same overall reweighting mechanism. What changes is how strongly individual hazard cues reduce the safety of movement actions and how much relative safety advantage is given to \texttt{stop}.
\\
\\
Both variants use the same traffic shield template. They differ only in how strongly the computed hazard signals influence the action-level safety scores. In the conservative variant, hazard cues impose stronger penalties on movement actions and give a larger relative advantage to \texttt{stop}. In the aggressive variant, the same hazard signals are scaled down through lower weighting, the safety gap among the movement actions is compressed more strongly, and the relative safety advantage of \texttt{stop} in benign states is reduced. In implementation terms, this is realized through adjusted weighting parameters such as the rule-weight map, the move-score blending strength, and the stop-margin terms. The aggressive shield therefore does not redefine what counts as hazardous; it only softens how strongly those hazards reshape the action distribution.

\subsection{Decentralized and Centralized Shields}
The thesis considers two shielding organizations. In the decentralized setting, each RL-controlled car is shielded separately from its own local symbolic observation, yielding a per-agent shielded distribution
\[
\pi_{\theta,i}^+(\cdot\mid o^i).
\]
\noindent This is consistent with decentralized execution and is modular to reuse across agents, since each car reasons only over locally visible traffic facts such as nearby pedestrians, cars, intersection occupancy, and priority cues.
\\
\\
The main limitation of decentralized shielding is coordination. Because one car's shield does not directly control the decisions of the other RL-controlled cars, symmetric intersection situations can lead to simultaneous hesitation even when a coordinated one-by-one passage would be safe.
\\
\\
In the centralized setting, the shield instead reasons over the RL-controlled cars jointly and produces a coordinated intervention at the multi-agent level. The motivation is not improved local hazard detection itself, but joint conflict resolution at shared intersections. Methodologically, centralized shielding preserves the same basic traffic-safety semantics as the decentralized template while extending it with coordinated passage decisions among the RL-controlled cars. This makes it the natural comparison point for testing whether direct joint coordination improves over decentralized probabilistic shielding in multi-agent autonomous driving.

\section{Reinforcement Learning Setup}
This section specifies how the multi-agent RL control problem is parameterized at the policy level. Since the thesis compares different multi-agent training organizations, the focus here is not on committing to one specific choice, but on formalizing the two training schemes considered for the LogiCity car population. In both cases, each controllable car acts from its own local grounded observation, selects one action from the common discrete action space, and receives task feedback from the shared traffic environment.

\subsection{Overview of the RL Loop}
Before distinguishing the two multi-agent training schemes, it is useful to summarize the common RL loop at a structural level. In both cases, the controllable cars act in the same shared LogiCity environment, receive local grounded observations, pass these observations through a policy-and-shield decision pipeline, execute the resulting actions simultaneously, and then collect trajectories that are used for policy optimization. The difference between the two training schemes lies not in the environment interaction itself, but in how decision modules and parameter updates are organized across the car population.
\\
\\
The shared-policy and separate-policy versions of this loop are compared below. The essential difference is whether all controllable cars feed into one common policy-update pipeline or whether each car maintains its own separate policy-update path.

\subsection{Shared-Policy and Separate-Policy Multi-Agent Training}
In the shared-policy setting, all controllable cars use one common stochastic policy
\[
\pi_\theta(a\mid o),
\]
\noindent so that each car samples its action as
\[
a_t^i \sim \pi_\theta(\cdot \mid o_t^i),
\qquad i\in\{1,\dots,K\}.
\]
\noindent The parameters \(\theta\) are therefore reused across the entire controllable car population, so all cars contribute experience to one shared update process. In the present LogiCity setting, this is natural because all controllable cars have the same action space, observation design, and task semantics.

\begin{figure}[H]
    \centering
    \includegraphics[width=0.98\textwidth]{figures/shared_policy_loop.png}
    \caption{Shared-policy multi-agent RL loop.}
    \label{fig:shared-policy-loop}
\end{figure}

\noindent In the separate-policy setting, each controllable car is assigned its own policy parameters. Instead of one shared controller, the system contains \(K\) distinct policies
\[
\pi_{\theta_1}(a\mid o),\;
\pi_{\theta_2}(a\mid o),\;
\dots,\;
\pi_{\theta_K}(a\mid o),
\]
\noindent and each car acts according to its own policy,
\[
a_t^i \sim \pi_{\theta_i}(\cdot \mid o_t^i),
\qquad i\in\{1,\dots,K\}.
\]
\noindent Here, the trajectories of car \(i\) are used to update \(\theta_i\) rather than a globally shared parameter set. This allows the car population to learn distinct control laws, but it also removes the statistical efficiency of pooled experience.

\begin{figure}[H]
    \centering
    \includegraphics[width=0.98\textwidth]{figures/separate_policy_loop.png}
    \caption{Separate-policy multi-agent RL loop.}
    \label{fig:separate-policy-loop}
\end{figure}

\noindent Methodologically, the distinction is therefore simple: shared-policy training enforces one common behavioral law across all controllable cars, whereas separate-policy training allows the car population to specialize. In the present thesis, both are evaluated under the same environment, observation design, and action space.

\section{Evaluation Protocol}
\subsection{Validation, Test, and Timeout Protocol}
Training is carried out online through repeated environment interaction, whereas validation and testing are performed on fixed pre-generated episode sets so that all compared methods are evaluated on the same traffic scenes. Validation is used for model selection; the final reported comparisons are based on the held-out test split. Episodes can end in success, failure, or timeout. The timeout threshold is normalized by an oracle trajectory length for the same scene:
\[
T_{\max}=2\,T_{\text{oracle}}
\]
\noindent so that long and short episodes are judged relative to their own difficulty rather than under one global horizon.
\\
\\
All quantitative comparisons are reported across different values of \(K\in\{1,2,3\}\), the number of RL-controlled cars in the environment. Across these settings, each scene always contains two rule-based pedestrians, while \(K\) controls only the number of RL cars.

\begin{table}[H]
\centering
\small
\caption{Evaluation metrics used in the thesis.}
\label{tab:evaluation-metrics}
\begin{tabularx}{\linewidth}{@{}l l X@{}}
\toprule
\textbf{Metric} & \textbf{Type} & \textbf{Role} \\
\midrule
Trajectory success rate (TSR) & Quantitative & Fraction of evaluated episodes in which all RL-controlled cars complete their routes successfully. Main scene-level performance metric. \\
\addlinespace
Failure rate & Quantitative & Fraction of episodes ending in safety failure, such as collision-related terminal events. \\
\addlinespace
Timeout rate & Quantitative & Fraction of episodes that do not fail explicitly but also do not finish within the oracle-scaled horizon. \\
\addlinespace
Action distribution & Shield-specific & Histogram over \texttt{slow}, \texttt{normal}, \texttt{fast}, and \texttt{stop}, used to characterize the behavioral effect of different shield configurations. \\
\bottomrule
\end{tabularx}
\end{table}

\chapter{Results}
\label{cha:results}

This chapter reports the empirical findings of the thesis and is organized around the four main comparisons studied in the experimental program. First, PLS is compared with unshielded RL as the number of simultaneously active RL-controlled cars increases. Second, decentralized shield variants are compared in order to study the trade-off between direct failure and timeout. Third, centralized shielding is analyzed on representative decentralized deadlock cases. Fourth, shared-policy and separate-policy training are compared under the same shielded traffic-control setting. Across the chapter, trajectory success rate remains the main task-completion metric, but it is always interpreted together with timeout behavior and failure modes.

\section{PLS Under Increasing Multi-Agent Interaction}
\label{sec:results-rq1}

The first results block addresses the most basic thesis question: whether PLS improves safe autonomous navigation relative to unshielded RL as multi-agent interaction pressure increases. The comparison is carried out across \(K=1\), \(K=2\), and \(K=3\), where larger \(K\) means that more RL-controlled cars must coordinate in the same shared traffic scene.

\subsection{Main Results}
\begin{table}[H]
\centering
\caption{Unshielded RL vs.\ PLS across increasing \(K\).}
\label{tab:results-rq1-main}
\begin{tabular}{llccc}
\toprule
\textbf{\(K\)} & \textbf{Method} & \textbf{TSR} & \textbf{Timeout rate} & \textbf{Failure rate} \\
\midrule
\(1\) & Unshielded PPO & 0.66 & 0.00 & 0.34 \\
\(1\) & PLS (Cons.) & 1.00 & 0.00 & 0.00 \\
\cmidrule(lr){1-5}
\(2\) & Unshielded PPO & 0.62 & 0.00 & 0.38 \\
\(2\) & PLS (Cons.) & 0.72 & 0.28 & 0.00 \\
\cmidrule(lr){1-5}
\(3\) & Unshielded PPO & 0.24 & 0.00 & 0.76 \\
\(3\) & PLS (Cons.) & 0.42 & 0.42 & 0.16 \\
\bottomrule
\end{tabular}
\end{table}

The fixed-test results show that conservative decentralized PLS improves both safety and task completion relative to unshielded PPO across all three evaluated values of \(K\). At \(K=1\), TSR rises from \(0.66\) to \(1.00\), while failure drops from \(0.34\) to \(0.00\). In the single-agent case, shielding therefore eliminates failure entirely and achieves perfect completion on the test split.
\\
\\
For larger \(K\), the same pattern remains visible, but the cost of conservatism becomes more pronounced. At \(K=2\), conservative PLS raises TSR from \(0.62\) to \(0.72\) and reduces failure from \(0.38\) to \(0.00\), although this comes with a timeout rate of \(0.28\). At \(K=3\), conservative PLS again improves over PPO, increasing TSR from \(0.24\) to \(0.42\) and reducing failure from \(0.76\) to \(0.16\), but now with timeout rate \(0.42\). The main conclusion is therefore that PLS is beneficial throughout the evaluated range of multi-agent interaction pressure, but that its advantage is increasingly mediated through a trade-off: direct unsafe failure is reduced strongly, while timeout becomes the main remaining limitation as the coordination problem becomes denser.

\section{Decentralized Shield Variants}
\label{sec:results-rq2}

The second results block studies how different decentralized shield variants affect multi-agent behavior. Here the central question is no longer whether shielding helps in general, but how different shield styles reshape the balance between safety preservation, efficiency, and coordination.
\\
\\
The comparison is carried out across \(K=1\), \(K=2\), and \(K=3\), so that the effect of conservative, aggressive, and mixed shield assignments can be interpreted under increasing multi-agent interaction pressure.

\subsection{Main Results}
\begin{table}[H]
\centering
\caption{Decentralized shield variants across increasing \(K\).}
\label{tab:results-rq2-main}
\begin{tabular}{llccc}
\toprule
\textbf{\(K\)} & \textbf{Shield setting} & \textbf{TSR} & \textbf{Timeout rate} & \textbf{Failure rate} \\
\midrule
\(1\) & Conservative & 1.00 & 0.00 & 0.00 \\
\(1\) & Aggressive & 0.66 & 0.00 & 0.34 \\
\cmidrule(lr){1-5}
\(2\) & Conservative & 0.72 & 0.28 & 0.00 \\
\(2\) & Aggressive & 0.64 & 0.00 & 0.36 \\
\(2\) & Mixed & 0.80 & 0.10 & 0.10 \\
\cmidrule(lr){1-5}
\(3\) & Conservative & 0.42 & 0.42 & 0.16 \\
\(3\) & Aggressive & 0.28 & 0.02 & 0.70 \\
\bottomrule
\end{tabular}
\end{table}

The decentralized shield comparison shows a clearer pattern than before: the conservative shield outperforms the aggressive shield in all reported settings. At \(K=1\), conservative shielding reaches perfect TSR and zero failure, whereas the aggressive variant remains at \(0.66\) TSR with \(0.34\) failure. At \(K=2\), conservative shielding again dominates aggressive shielding, improving TSR from \(0.64\) to \(0.72\) and reducing failure from \(0.36\) to \(0.00\), although with a non-negligible timeout rate. At \(K=3\), the same ordering remains visible: the conservative shield achieves higher TSR and much lower failure than the aggressive one. The main distinction is therefore not that aggressive shielding is preferable at some values of \(K\), but that conservative shielding is consistently stronger overall while paying for this with more waiting behavior.
\\
\\
The most informative case is \(K=2\), where the mixed setting improves further over both fully conservative and fully aggressive shielding, reaching TSR \(0.80\) with timeout rate \(0.10\) and failure rate \(0.10\). A plausible explanation is that the mixed assignment partially breaks the behavioral symmetry between the two RL cars. Under fully symmetric decentralized shielding, both cars are more likely to react in the same way at the same time, which can reinforce mutual hesitation and increase deadlock-like timeout behavior. The mixed configuration weakens this symmetry by making the two cars respond somewhat differently to the same local traffic situation, which can help create staggered timing and allow one agent to progress while the other yields. In this sense, the \(K=2\) result suggests that decentralized multi-agent shielding is influenced not only by how conservative the shield is in isolation, but also by whether the resulting interaction dynamics amplify or break symmetric waiting patterns.

\section{Centralized Shielding in Deadlock Cases}
\label{sec:results-rq3}

The third results block focuses on the coordination cases in which decentralized shielding becomes most challenged. In particular, the analysis considers representative multi-agent intersection episodes in which decentralized shielding exhibits deadlock-like hesitation or timeout behavior. These episodes are especially informative because they isolate the situations in which local per-agent shielding is no longer sufficient and joint coordination becomes the central issue.
\\
\\
The purpose of this analysis is to study whether centralized shielding can resolve such coordination failures by treating the RL-controlled cars jointly. The main question is therefore whether centralized shielding changes the \emph{type} of multi-agent behavior by enabling ordered joint conflict resolution in cases where decentralized shielding leaves the agents in mutual waiting.

\subsection{Targeted Results}
\begin{table}[H]
\centering
\caption{Centralized shielding on representative decentralized deadlock cases.}
\label{tab:results-rq3-main}
\begin{tabular}{lcc}
\toprule
\textbf{Episode} & \textbf{Decentralized outcome} & \textbf{Centralized outcome} \\
\midrule
\texttt{ep0} & timeout / deadlock & success \\
\texttt{ep9} & timeout / deadlock & success \\
\texttt{ep10} & timeout / deadlock & success \\
\texttt{ep13} & timeout / deadlock & success \\
\bottomrule
\end{tabular}
\end{table}

\begin{figure}[H]
    \centering
    \includegraphics[width=0.50\linewidth]{figures/results_centralized_ep0_step48.png}
    \caption{Representative coordination state from \texttt{ep0} at step 48.}
    \label{fig:centralized-ep0-coordination-state}
\end{figure}

Across these selected cases, the centralized shield is evaluated precisely on the interaction situations for which joint coordination is most relevant. Episodes \texttt{ep0}, \texttt{ep9}, \texttt{ep10}, and \texttt{ep13} were chosen because decentralized shielding became stuck on all of them, whereas centralized shielding completed them successfully. The targeted comparison therefore isolates the coordination cases in which the difference between decentralized and centralized shielding is most visible.
\\
\\
Figure~\ref{fig:centralized-ep0-coordination-state} illustrates the key coordination situation underlying this result. In \texttt{ep0}, the traffic state reaches a moment at which all three RL-controlled cars are simultaneously visible in the same local interaction region. From the viewpoint of each decentralized agent, the local observation now contains the other two RL cars as active conflict partners at or near the same intersection, together with priority and occupancy cues indicating that entering is risky. The decentralized shield reacts exactly to these local conflict signals: for each car individually, it increases the probability of \texttt{stop} and suppresses forward motion. Since all three shields perform this same conservative reaction independently, no car is explicitly released to break the symmetry, and the interaction collapses into mutual waiting.
\\
\\
This does not mean that decentralized shielding always deadlocks. When the approach timing is slightly asymmetric, so that the cars arrive more sequentially, the local observations are no longer perfectly symmetric and decentralized shielding can still resolve the interaction successfully. The deadlock risk becomes especially pronounced as the number of RL-controlled cars increases and several of them reach the same intersection at nearly the same time, because then the local stop incentives become aligned across agents.
\\
\\
Under centralized shielding, by contrast, the shield receives this same three-car interaction structure but does not evaluate it as three separate stop decisions. Instead, it treats the configuration as one joint coordination problem and can hold conflicting cars while releasing the currently preferred car through the intersection. In this way, the centralized shield still reacts conservatively to the shared hazard, but it converts that caution into an ordered one-by-one passage rather than a symmetric deadlock. The important result is therefore not merely that centralized shielding attains a better final outcome label on these episodes, but that it changes the coordination mechanism itself. This targeted analysis shows directly how centralized shielding can resolve a class of decentralized multi-agent deadlocks that arises when several RL-controlled cars meet at the same intersection at the same time.

\section{Shared vs.\ Separate Policy Training}
\label{sec:results-rq4}

The fourth results block compares shared-policy and separate-policy multi-agent training under PLS. The central question is whether the organization of learning changes the resulting multi-agent behavior, even when the task, shield, and environment remain the same.
\\
\\
This comparison also allows the thesis to study whether separately trained agents actually diverge behaviorally from one another or remain effectively interchangeable despite being optimized independently.

\subsection{Main Results}
\begin{table}[H]
\centering
\caption{Shared-policy vs.\ separate-policy training under PLS.}
\label{tab:results-rq4-main}
\begin{tabular}{lccc}
\toprule
\textbf{Training scheme} & \textbf{TSR} & \textbf{Timeout rate} & \textbf{Failure rate} \\
\midrule
Shared-policy PLS & 0.42 & 0.42 & 0.16 \\
Separate-policy PLS & 0.14 & 0.46 & 0.40 \\
\bottomrule
\end{tabular}
\end{table}

For the shared-policy reference, the \(K=3\) conservative decentralized shield result is used, since this is the strongest shared-policy configuration available in the current multi-agent benchmark. The separate-policy condition is evaluated on the fixed test split with the canonical identity assignment, that is, each RL-controlled car is paired with its own trained policy rather than with permuted policy assignments.
\\
\\
Under this fixed-test comparison, shared-policy PLS reaches a TSR of \(0.42\), with timeout rate \(0.42\) and failure rate \(0.16\). Separate-policy PLS performs substantially worse, reaching TSR \(0.14\), timeout rate \(0.46\), and failure rate \(0.40\). At the aggregate level, the shared-policy setup therefore yields clearly more stable multi-agent coordination and substantially better overall task completion.

\begin{figure}[H]
    \centering
    \includegraphics[width=0.82\linewidth]{figures/results_separate_policy_action_hist.png}
    \caption{Per-policy action distributions for separate-policy \(K=3\).}
    \label{fig:results-separate-policy-action-hist}
\end{figure}

\noindent This result is consistent with the structure of the task itself. In the present benchmark, all RL-controlled cars face the same control problem: they share the same action space, the same observation design, the same reward structure, and the same general traffic objective. In such a setting, parameter sharing is not a restriction but a natural inductive bias. One common policy is sufficient to represent the required behavior, and pooling experience across all cars also makes learning more data-efficient and more stable. Separate policies would be more compelling in a setting with genuinely different agent roles, capabilities, or observation regimes. Here, however, the agents are intentionally homogeneous, so forcing them to learn independently mainly removes useful regularization.
\\
\\
The action distributions support exactly this interpretation. Figure~\ref{fig:results-separate-policy-action-hist} shows that separate training breaks behavioral symmetry strongly: the three policies develop visibly different action profiles, with one policy becoming much more stop-dominated and another shifting toward a slow-stop pattern. This shows that separate training does create specialization, but in the current task that specialization is not beneficial. Instead, the learned behaviors become less aligned with one another, which makes joint coordination less stable and reduces scene-level performance. The main conclusion is therefore that, for homogeneous multi-agent traffic control of this kind, shared-policy training is not only simpler but also the more effective organization of learning.

\section{Summary}
\label{sec:results-summary}

The results of this chapter show that probabilistic logic shielding improves safe multi-agent navigation in LogiCity, but that its effect is best understood as a coordination trade-off rather than as a purely safety-only gain. Relative to unshielded PPO, conservative decentralized PLS consistently reduced failure and improved trajectory success, although at higher values of \(K\) part of this gain was offset by increased timeout behavior.
\\
\\
Among decentralized shield variants, the conservative shield performed better than the aggressive shield across all reported settings, while the mixed \(K=2\) assignment improved further by reducing symmetric hesitation between agents. The targeted centralized analysis then showed that some remaining failures of decentralized shielding are genuine coordination deadlocks, and that centralized shielding can resolve such cases by enforcing coordinated one-by-one passage through shared intersections.
\\
\\
Finally, the comparison between shared-policy and separate-policy training showed that homogeneous multi-agent traffic control is better served by parameter sharing than by separate specialization. Although separate policies developed distinct behaviors, this divergence reduced coordination quality and lowered overall scene-level performance. Taken together, the chapter therefore supports the central claim of the thesis: PLS is an effective mechanism for safe multi-agent traffic control, but its practical value depends on how shielding interacts with coordination structure, symmetry, and learning organization.

\chapter{Conclusion}
\label{cha:conclusion}

This thesis studied Probabilistic Logic Shielding (PLS) for safe multi-agent autonomous driving in LogiCity. The focus was not only on reducing unsafe behavior, but on understanding how shielding affects \emph{safe task completion} when multiple RL-controlled cars must coordinate under shared traffic rules. The experiments were organized around four main questions: how PLS compares with unshielded RL as the number of RL-controlled cars increases, how decentralized shield variants trade off safety against efficiency, whether centralized shielding can resolve coordination failures that remain difficult under decentralized shielding, and how shared-policy and separate-policy training differ under PLS.
\\
\\
The results show that PLS changes multi-agent behavior in a meaningful way, but not in a uniformly one-dimensional manner. Relative to unshielded PPO, decentralized conservative PLS substantially reduced direct failures and improved joint trajectory success at \(K=1\) and \(K=3\). At the same time, it also introduced a clear conservatism cost: as the number of RL-controlled cars increased, part of the unsafe failure mass was shifted into timeout behavior. This means that the main role of PLS in the current benchmark is not simply to ``make the policy safer,'' but to reshape the balance between unsafe failure and cautious incompletion.
\\
\\
The comparison of decentralized shield variants confirmed this interpretation. No single decentralized shield style dominated across all settings. Conservative shielding reduced failures most strongly but became increasingly vulnerable to coordination-induced timeouts, whereas aggressive shielding preserved movement better but at the cost of substantially more unsafe outcomes. The mixed \(K=2\) setting was particularly informative: it outperformed both fully conservative and fully aggressive decentralized shielding, suggesting that part of the difficulty is not only the strength of the shield itself, but also whether the resulting interaction dynamics reinforce or break symmetric waiting behavior.
\\
\\
The targeted centralized analysis showed that some of these failures are fundamentally coordination failures rather than failures of local hazard detection. In representative episode-level comparisons, decentralized shielding became stuck when several RL-controlled cars reached the same intersection in near-symmetric timing, because each local shield independently increased stopping pressure without explicitly resolving which vehicle should proceed. Centralized shielding, by contrast, could treat the same interaction as a joint coordination problem and impose an ordered one-by-one passage. In this sense, centralized shielding changed not only the final outcome but the coordination mechanism itself.
\\
\\
The comparison between shared-policy and separate-policy training led to a similarly structural conclusion. In the present benchmark, all RL-controlled cars have the same task, observation design, and action space. Under these conditions, shared-policy training proved more effective than separate-policy training. The separate-policy case did produce clear behavioral divergence across the three learned policies, but this specialization did not improve collective performance. Instead, it reduced coordination stability and lowered scene-level trajectory success. For homogeneous multi-agent traffic control of this kind, parameter sharing therefore appears to be the more suitable learning organization.
\\
\\
Overall, the thesis shows that the value of probabilistic logic shielding in multi-agent urban navigation cannot be judged through failure reduction alone. In logic-driven traffic environments, shielding must be evaluated through the combined lens of success, failure, timeout, and coordination. The main conclusion is therefore that PLS is a useful and promising safety mechanism for multi-agent autonomous driving, but that its effectiveness depends strongly on how shielding interacts with multi-agent symmetry, coordination structure, and policy organization.

\section{Limitations}

Several limitations constrain the current study. First, the experiments were conducted in a stylized compact LogiCity benchmark centered around intersection scenarios. This setting is suitable for controlled analysis of multi-agent interaction, but it remains simpler than real urban traffic with richer maps, traffic lights, additional semantic vehicle classes, and more diverse interaction patterns.
\\
\\
Second, the main decentralized shielding design relies on local symbolic observations. This is a natural and scalable design choice, but it also creates a coordination limitation: when several RL-controlled cars encounter the same interaction structure with near-symmetric timing, their local shields may react in similarly conservative ways without any explicit mechanism for joint release. The deadlock cases studied in the centralized analysis illustrate this limitation directly.
\\
\\
Third, although the targeted centralized analysis shows that centralized shielding can resolve some decentralized deadlocks, centralized shielding also introduces its own practical disadvantages. It is more complex to design, computationally heavier at runtime, and more difficult to scale robustly as the number of interacting agents and edge cases grows. By contrast, decentralized shielding is easier to specify, modular to reuse across agents, and computationally simpler. The current thesis therefore does not treat centralized shielding as a universally superior replacement, but rather as a coordination-oriented alternative with its own costs.
\\
\\
Fourth, although the thesis includes an uncertain symbolic observation model, it still operates within a simulator in which the symbolic abstraction and traffic-rule structure are known and controlled. This is valuable for studying interpretable shielding, but it remains a step away from full real-world perception pipelines and deployment conditions.

\section{Future Work}

Several extensions follow naturally from these results. A first direction is to investigate richer forms of multi-agent shielding. The current results suggest that decentralized local shielding is effective for many situations but struggles in symmetric coordination cases, while centralized shielding can resolve some of these deadlocks but at higher design and computational cost. Future work could therefore study stronger centralized formulations, hybrid decentralized-centralized schemes, or factored regional shields as discussed in prior multi-agent shielding work \cite{elsayedaly2021safemultiagentreinforcementlearning}.
\\
\\
A second direction is to extend the current benchmark toward more realistic traffic structure. This includes larger maps, denser deployments, more traffic-participant types, and additional traffic-control elements. Since LogiCity supports richer semantic abstractions \cite{li2025logicityadvancingneurosymbolicai}, it provides a natural platform for such scaling studies.
\\
\\
A third direction is to study more diverse multi-agent learning organizations. The current results indicate that shared-policy training is well suited to homogeneous traffic agents, whereas separate-policy specialization did not help in the present task. Future work could revisit this question in settings with genuinely different agent roles, capabilities, or observation regimes, where specialization may become more useful.
\\
\\
Finally, a promising longer-term direction is to bring probabilistic logic shielding closer to end-to-end autonomous driving pipelines. This includes tighter integration with perception uncertainty, richer symbolic sensor extraction, and stronger coupling between learned policies and structured safety reasoning. Such work would help clarify how far the advantages of PLS in interpretable simulated traffic transfer to more realistic autonomous driving systems.


% If you have appendices:
\appendixpage*          % if wanted
\appendix
\chapter{The First Appendix}
\label{app:A}
% Appendices hold useful data which is not essential to understand the work
% done in the master's thesis. An example is a (program) source.
% An appendix can also have sections as well as figures and references\cite{h2g2}.

\section{Predicate Definitions}
\label{app:predicate-definitions}

\begin{table}[H]
\centering
\small
\caption{Predicate definitions used in the current thesis experiments.}
\label{tab:predicate-definitions}
\begin{tabularx}{\linewidth}{@{}p{3.8cm} X@{}}
\toprule
\textbf{Predicate} & \textbf{Meaning} \\
\midrule

\texttt{IsPedestrian(x)}
& True iff entity \texttt{x} is a pedestrian. \\

\texttt{IsCar(x)}
& True iff entity \texttt{x} is a car. \\

\texttt{IsAtInter(x)}
& True iff agent \texttt{x} is at the entrance or boundary of an intersection, i.e. in the approach region before fully entering it. \\

\texttt{IsInInter(x)}
& True iff agent \texttt{x} is inside the intersection region itself. \\

\texttt{IsParked(x)}
& True iff car \texttt{x} is inactive because it is still parked at spawn or already parked after reaching its goal. This predicate is used internally to exclude parked cars from active conflict and failure evaluation. \\

\texttt{SameInter(x,y)}
& True iff agents \texttt{x} and \texttt{y} are associated with the same intersection region. This predicate is part of the observation and is used by the traffic logic to restrict conflict reasoning to the relevant intersection. \\

\texttt{HigherPri(x,y)}
& True iff \texttt{x} has higher traffic priority than \texttt{y}. In the implementation, this corresponds to \texttt{x} having a numerically smaller priority value than \texttt{y}. \\

\texttt{CollidingClose(x,y)}
& True iff agent \texttt{x} is in an imminent collision-close relation with agent \texttt{y}, meaning that continuing movement would constitute an immediate hazard cue. \\

\texttt{Collision(x,y)}
& True iff agent \texttt{x} is in a realized collision state with agent \texttt{y}. This predicate is used for terminal failure detection. \\

\bottomrule
\end{tabularx}
\end{table}

\section{Observation Noise Model}
\label{app:observation-noise-model}

This appendix specifies the observation-noise model used in the thesis. The purpose of this model is to introduce uncertainty into the \emph{environment-side symbolic observation} received by the agents, thereby mimicking imperfect sensing in autonomous driving. It is not a component of PLS itself. Instead, the environment first constructs a grounded symbolic observation, after which the observation-noise model corrupts that grounded vector before it is passed to the policy and, in shielded runs, also to the shield.

\subsection{Placement in the Pipeline}

Let \(o \in [0,1]^d\) denote the grounded symbolic observation vector before corruption. The observation-noise model is applied after predicate grounding and before policy evaluation. Thus, the simulator dynamics, map structure, route generation, and logical task rules remain unchanged. Only the symbolic observation channel is made uncertain.
\\
\\
The observation is transformed into a confidence-valued symbolic vector. In this way, the grounded logical facts remain directly interpretable, but they are no longer perfectly certain.

\subsection{Entity-Correlated Asymmetric Construction}

The implementation reconstructs entity-wise groups from the predicate-grounding layout. Unary grounded entries are assigned to per-entity type or relation buckets, and binary grounded entries are assigned to both participating entities. This induces a set of grouped symbolic facts for the ego entity and for each non-ego entity in the local field of view.
\\
\\
The model distinguishes between:
\begin{itemize}
    \item \textbf{ego facts},
    \item \textbf{non-ego type facts},
    \item \textbf{non-ego relation facts}.
\end{itemize}

\noindent It is asymmetric in two senses:
\begin{itemize}
    \item false negatives and false positives use different rates,
    \item type predicates and relation predicates use different rates.
\end{itemize}

\subsection{Confidence Mapping}

The model does not directly flip each fact into a new binary alternative. Instead, it maps each grounded fact to a confidence value.
\\
\\
For a given fact \(x \in \{0,1\}\), the confidence-valued output \(\tilde{x}\) is constructed as
\[
\tilde{x} =
\begin{cases}
1 - p_{\mathrm{fn}}, & x = 1,\\
p_{\mathrm{fp}}, & x = 0,
\end{cases}
\]
where \(p_{\mathrm{fn}}\) and \(p_{\mathrm{fp}}\) depend on the fact category and, for non-ego relation facts, also on whether the corresponding entity is treated as occluded.
\\
\\
Thus, true facts become high-confidence values below \(1\), and false facts become low-confidence values above \(0\). The output remains numeric and can therefore be consumed directly by both the policy and the shield.

\subsection{Occlusion and Entity Dropout}

For every non-ego entity, the implementation samples an occlusion event:
\[
z_i \sim \mathrm{Bernoulli}(p_i^{\mathrm{drop}}),
\]
where \(z_i = 1\) indicates that entity \(i\) is treated as occluded. The effective dropout probability is not fixed globally; it increases with local crowding. Let \(n_{\mathrm{non\text{-}ego}}\) denote the number of present non-ego entities in the current observation and let \(n_{\max}\) denote the maximum possible number of non-ego entities under the observation budget. The implementation computes a crowding ratio
\[
r_{\mathrm{crowd}} =
\begin{cases}
0, & n_{\mathrm{non\text{-}ego}} \le 1,\\
\dfrac{n_{\mathrm{non\text{-}ego}} - 1}{\max(n_{\max} - 1, 1)}, & n_{\mathrm{non\text{-}ego}} > 1.
\end{cases}
\]
Using this ratio, the effective entity-dropout probability becomes
\[
p_i^{\mathrm{drop}}

=
\mathrm{clip}\!\left(
p_{\mathrm{drop}} \cdot
\left(1 + g_{\mathrm{drop}}\, r_{\mathrm{crowd}}\right),
0, 1
\right),
\]
where \(p_{\mathrm{drop}}\) is the base entity-dropout parameter and \(g_{\mathrm{drop}}\) is the crowding-related dropout gain.

\subsection{Crowding-Dependent Relation Noise}

The model also increases relation uncertainty when more non-ego entities populate the observation. A crowding-dependent relation-noise scale is computed as
\[
\lambda_{\mathrm{rel}} = 1 + g_{\mathrm{rel}}\, r_{\mathrm{crowd}},
\]
where \(g_{\mathrm{rel}}\) is the crowding-related relation-noise gain.
\\
\\
For non-ego relation facts, the active false-negative and false-positive rates are then
\[
p_{\mathrm{fn}}^{\mathrm{rel},i} =
\mathrm{clip}\!\left(
\bar{p}_{\mathrm{fn}}^{\mathrm{rel},i}\,\lambda_{\mathrm{rel}},
0, 1
\right),
\qquad
p_{\mathrm{fp}}^{\mathrm{rel},i} =
\mathrm{clip}\!\left(
\bar{p}_{\mathrm{fp}}^{\mathrm{rel},i}\,\lambda_{\mathrm{rel}},
0, 1
\right),
\]
where
\[
\bar{p}_{\mathrm{fn}}^{\mathrm{rel},i}
=
\begin{cases}
p_{\mathrm{fn}}^{\mathrm{rel}}, & z_i = 0,\\
p_{\mathrm{fn}}^{\mathrm{rel,occ}}, & z_i = 1,
\end{cases}
\qquad
\bar{p}_{\mathrm{fp}}^{\mathrm{rel},i}
=
\begin{cases}
p_{\mathrm{fp}}^{\mathrm{rel}}, & z_i = 0,\\
p_{\mathrm{fp}}^{\mathrm{rel,occ}}, & z_i = 1.
\end{cases}
\]
\noindent Hence, occluded entities receive more aggressive relation corruption, and this effect becomes even stronger in more crowded local observations.

\subsection{Predicate Categories}

In the experiments of this thesis, the relevant \textbf{type predicates} are
\begin{center}
\texttt{IsCar}, \texttt{IsPedestrian}.
\end{center}

\noindent The relevant \textbf{relation and traffic-state predicates} are the grounded traffic facts used by the cars and the shield, in particular predicates such as
\begin{center}
\texttt{IsAtInter}, \texttt{IsInInter}, \texttt{HigherPri}, \texttt{CollidingClose},
\end{center}
\noindent together with the same-intersection relation used by the traffic logic. This distinction matters because the most safety-critical facts in the thesis belong to the relation category and are therefore perturbed more strongly than simple semantic identity facts.

\subsection{Parameter Values Used in the Thesis}

The current thesis experiments use the parameterization shown in Table~\ref{tab:observation-noise-parameters}. These values are shared across the main experiment configurations.

\begin{table}[H]
\centering
\small
\caption{Observation-noise parameters used in the thesis experiments.}
\label{tab:observation-noise-parameters}
\begin{tabularx}{\linewidth}{@{}l c X@{}}
\toprule
\textbf{Parameter} & \textbf{Value} & \textbf{Role} \\
\midrule
\texttt{entity\_dropout} & \(0.15\) & Base occlusion/dropout probability for each non-ego entity. \\
\texttt{false\_negative\_type} & \(0.03\) & False-negative rate for non-ego type predicates. \\
\texttt{false\_positive\_type} & \(0.01\) & False-positive rate for non-ego type predicates. \\
\texttt{false\_negative\_relation} & \(0.10\) & Base false-negative rate for non-occluded non-ego relation predicates. \\
\texttt{false\_positive\_relation} & \(0.03\) & Base false-positive rate for non-occluded non-ego relation predicates. \\
\texttt{occluded\_false\_negative\_relation} & \(0.20\) & Base false-negative rate for occluded non-ego relation predicates. \\
\texttt{occluded\_false\_positive\_relation} & \(0.05\) & Base false-positive rate for occluded non-ego relation predicates. \\
\texttt{crowding\_relation\_noise\_gain} & \(0.75\) & Strength of the increase in relation uncertainty under higher local crowding. \\
\texttt{crowding\_entity\_dropout\_gain} & \(0.50\) & Strength of the increase in non-ego dropout under higher local crowding. \\
\texttt{ego\_false\_negative} & \(0.01\) & False-negative rate for ego-related facts. \\
\texttt{ego\_false\_positive} & \(0.00\) & False-positive rate for ego-related facts. \\
\bottomrule
\end{tabularx}
\end{table}

\subsection{Interpretation}

The resulting observation model is designed to reflect a simple but structured intuition about autonomous driving perception. Semantic identity facts are comparatively stable, while relational traffic facts are more fragile. Ego-related facts are the most reliable. Non-ego relation facts become less reliable under occlusion, and local crowding increases both the chance that a non-ego entity is effectively dropped and the uncertainty attached to its relational facts.
\\
\\
The purpose of this model is therefore not to simulate a full sensor stack, but to provide a controlled symbolic uncertainty process that is rich enough to stress-test PLS in multi-agent traffic scenarios while remaining interpretable and directly linked to the grounded logical observation space.

\section{PLPG Safety Regularizer}
\label{app:safety-regularizer}

This appendix specifies the safety-regularization term used in the current PLPG implementation. The purpose of this term is to bias optimization toward policy distributions that are safer under the current symbolic traffic facts, beyond the effect already produced by return maximization alone.

\subsection{Safety Quantity Used in Training}

For each state \(s\), the shield first computes action-level safety scores \(q(s,a)\) for all actions \(a \in \mathcal{A}\). From these scores, two policy-level safety quantities can be formed:
\[
P_{\pi_\theta}(\texttt{safe}\mid s)
=
\sum_{a \in \mathcal{A}} \pi_\theta(a\mid s)\, q(s,a),
\]
\[
P_{\pi_\theta^+}(\texttt{safe}\mid s)
=
\sum_{a \in \mathcal{A}} \pi_\theta^+(a\mid s)\, q(s,a).
\]
\noindent The first quantity measures the expected safety of the raw base policy, whereas the second measures the expected safety of the shielded policy.

\subsection{Implemented Loss Term}

In the current thesis implementation, the safety regularizer is applied as a negative log-safety penalty:
\[
\mathcal{L}_{\mathrm{safety}}
=
-\log\!\big(P(\texttt{safe}\mid s)\big).
\]
\noindent Averaged over a minibatch, this becomes
\[
\mathcal{L}_{\mathrm{safety}}^{\mathrm{batch}}
=
\frac{1}{B}\sum_{t=1}^{B}
-\log\!\big(P(\texttt{safe}\mid s_t)\big).
\]
\noindent By default, the implementation uses the shielded policy safety
\[
P(\texttt{safe}\mid s_t)=P_{\pi_\theta^+}(\texttt{safe}\mid s_t),
\]
\noindent although the code also permits regularization with the base-policy safety quantity.

\subsection{Combined PLPG Objective}

The final training loss is the PPO objective augmented with the safety regularizer:
\[
\mathcal{L}
=
\mathcal{L}_{\mathrm{policy}}
+
c_v\,\mathcal{L}_{\mathrm{value}}
+
c_e\,\mathcal{L}_{\mathrm{entropy}}
+
\alpha\,\mathcal{L}_{\mathrm{safety}},
\]
\noindent where \(c_v\) and \(c_e\) are the usual value-loss and entropy-loss coefficients, and \(\alpha\) controls the strength of safety regularization.

\subsection{Interpretation}

This regularizer does not punish only the single sampled action. Instead, it penalizes the entire action distribution when its expected safety is low. Policies that place more probability mass on actions with low action-level safety scores receive a larger loss. The use of the negative logarithm makes this effect stronger when the policy becomes highly unsafe, while only mildly penalizing policies whose safety is already high. In this way, the safety term encourages the policy to shift probability mass toward safer actions at the distribution level.

\section{Implemented Soft-Traffic Hazard Calculation}
\label{app:soft-traffic-hazard-calculation}

This appendix specifies the standard decentralized \texttt{easy\_soft\_traffic} template. The calculation is performed independently for each ego car and each visible non-ego entity \(j\). All predicate values are confidences in \([0,1]\). Let \(c\), \(a\), and \(r\) be the ego confidences for \texttt{IsCar}, \texttt{IsAtInter}, and \texttt{IsInInter}. Let \(a_j,r_j,s_j,p_j,h_j,d_j\) be, respectively, the other entity's at-intersection, in-intersection, same-intersection, pedestrian, higher-priority, and collision-close confidences; \(h_j\) denotes \texttt{HigherPri(j, ego)}. Define \(\operatorname{clip}(x)=\min(\max(x,0),1)\).

\subsection{Per-Entity Conflict Cues}

The intersection-overlap cue is
\[
o_j=\max\{a\,a_j,a\,r_j,r\,a_j,r\,r_j\}.
\]
The template then calculates
\[
\begin{aligned}
d^{\mathrm{ped}}_j &= c\,p_j\,s_j\,o_j, &
d^{\mathrm{occ}}_j &= a\,s_j\,r_j(1-0.5p_j),\\
d^{\mathrm{pri}}_j &= a\,a_j\,s_j\,h_j(1-p_j), &
d^{\mathrm{sim}}_j &= c(1-p_j)s_j\,r\,r_j,\\
d^{\mathrm{pedtr}}_j &= c\,p_j\,s_j\max\{r_j,0.70d_j,0.35a_j\}, &
d^{\mathrm{vehtr}}_j &= c(1-p_j)s_j\max\{\max(r_j,a_j(1-p_j)),0.60d_j\}.
\end{aligned}
\]
These represent pedestrian, occupied-intersection, higher-priority, simultaneous-occupancy, pedestrian-transition, and vehicle-transition cues, respectively.

\subsection{Aggregated Hazards and Traffic Regimes}

The most severe visible entity determines each hazard:
\[
\begin{aligned}
H_{\mathrm{cc}}&=\operatorname{clip}(\max_j 0.85d_j), &
H_{\mathrm{ped}}&=\operatorname{clip}(\max_j d^{\mathrm{ped}}_j),\\
H_{\mathrm{occ}}&=\operatorname{clip}(\max_j d^{\mathrm{occ}}_j), &
H_{\mathrm{pri}}&=\operatorname{clip}(\max_j d^{\mathrm{pri}}_j),\\
H_{\mathrm{sim}}&=\operatorname{clip}(\max_j d^{\mathrm{sim}}_j), &
H_{\mathrm{tr}}&=\operatorname{clip}(\max_j\max\{d^{\mathrm{pedtr}}_j,d^{\mathrm{vehtr}}_j\}),\\
H&=\max\{H_{\mathrm{cc}},H_{\mathrm{ped}},H_{\mathrm{occ}},H_{\mathrm{pri}},H_{\mathrm{sim}},H_{\mathrm{tr}}\}.
\end{aligned}
\]
With \(U=\max\{H_{\mathrm{ped}},H_{\mathrm{occ}},H_{\mathrm{pri}},H_{\mathrm{sim}}\}\), the four regimes are
\[
\begin{aligned}
F &= \operatorname{clip}(1-\max\{U,H_{\mathrm{tr}},H_{\mathrm{cc}}\}),\\
C &= \operatorname{clip}\!\left(\max\{U,H_{\mathrm{tr}}\}\operatorname{clip}\!\left(\frac{H_{\mathrm{tr}}-0.12}{0.35}\right)\right),\\
E &= \operatorname{clip}\!\left(\max\{a,0.75r\}\max\{U,H_{\mathrm{tr}}\}\operatorname{clip}\!\left(\frac{H_{\mathrm{tr}}-0.28}{0.30}\right)\right),\\
K &= \operatorname{clip}\!\left(\max\left\{H_{\mathrm{cc}},H_{\mathrm{sim}},0.9a\max\{H_{\mathrm{ped}},H_{\mathrm{pri}},H_{\mathrm{occ}}\},\operatorname{clip}\!\left(\frac{H_{\mathrm{tr}}-0.52}{0.20}\right)\right\}\right).
\end{aligned}
\]
Here \(F,C,E,K\) denote free flow, caution, conflict edge, and critical traffic.

\subsection{Action-Level Safety Scores}

The movement-action base risks are
\[
\begin{aligned}
B_{\mathrm{slow}}&=\max\{0.28H_{\mathrm{cc}},0.12H_{\mathrm{pri}},0.16H_{\mathrm{occ}},0.24H_{\mathrm{ped}},0.52H_{\mathrm{sim}},0.34H_{\mathrm{tr}}\},\\
B_{\mathrm{normal}}&=\max\{0.72H_{\mathrm{cc}},0.42H_{\mathrm{pri}},0.52H_{\mathrm{occ}},0.64H_{\mathrm{ped}},0.86H_{\mathrm{sim}},0.82H_{\mathrm{tr}}\},\\
B_{\mathrm{fast}}&=\max\{0.96H_{\mathrm{cc}},0.88H_{\mathrm{pri}},0.94H_{\mathrm{occ}},0.98H_{\mathrm{ped}},H_{\mathrm{sim}},H_{\mathrm{tr}}\}.
\end{aligned}
\]
Regime shaping yields
\[
\begin{aligned}
R_{\mathrm{slow}}&=\max\{B_{\mathrm{slow}},0.03F,0.10C,0.24E,0.55K\},\\
R_{\mathrm{normal}}&=\max\{B_{\mathrm{normal}},0.04F,0.48C,0.82E,0.97K\},\\
R_{\mathrm{fast}}&=\max\{B_{\mathrm{fast}},0.06F,0.88C,0.95E,K\}.
\end{aligned}
\]
Finally, the exact scores used for reweighting are
\[
\begin{aligned}
q(o^i,\texttt{slow})&=\operatorname{clip}(1-R_{\mathrm{slow}}),\\
q(o^i,\texttt{normal})&=\operatorname{clip}(1-R_{\mathrm{normal}}),\\
q(o^i,\texttt{fast})&=\operatorname{clip}(1-R_{\mathrm{fast}}),\\
q(o^i,\texttt{stop})&=1.
\end{aligned}
\]
Thus \texttt{stop} always has the maximum raw score in the standard template. The increasingly larger movement coefficients ensure that \texttt{fast} loses safety before \texttt{normal}, while \texttt{slow} remains the safest movement action.

\backmatter
% The bibliography comes after the appendices.
% You can replace the standard "abbrv" bibliography style by another one.
\bibliographystyle{abbrv}
\bibliography{references}
\include{CoC}
\end{document}
